% Constructive Repair in Realization Calculus
% Standalone LaTeX monograph. Compile with pdfLaTeX (three passes).
% Author: Komorebi P. Mathematical source: Constructive Repair Lean v0.5.0, retaining v0.2.0--v0.4.0.
\documentclass[11pt,letterpaper,oneside,openany]{book}
\usepackage[T1]{fontenc}
\usepackage[utf8]{inputenc}
\usepackage{lmodern}
\usepackage[scaled=0.91]{helvet}
\usepackage{microtype}
\usepackage[margin=1in,headheight=15pt,headsep=19pt,footskip=29pt]{geometry}
\usepackage{amsmath,amssymb,amsthm,mathtools}
\usepackage{bm}
\usepackage[table]{xcolor}
\definecolor{Pine}{HTML}{244D43}
\definecolor{Ink}{HTML}{192D2B}
\definecolor{Slate}{HTML}{596965}
\definecolor{Ochre}{HTML}{92703E}
\definecolor{WarmPaper}{HTML}{F4F2EC}
\definecolor{PalePine}{HTML}{F0F5F2}
\definecolor{RuleGray}{HTML}{D8E0DB}
\usepackage{booktabs,array,tabularx,longtable}
\usepackage{enumitem}
\usepackage{titlesec}
\usepackage{fancyhdr}
\usepackage{needspace}
\usepackage{tikz-cd}
\usepackage[most]{tcolorbox}
\usepackage{fvextra}
\usepackage{xurl}
\usepackage[colorlinks=true,linkcolor=Pine,citecolor=Pine,urlcolor=Pine,
  bookmarksnumbered=true,bookmarksopen=false,pdfdisplaydoctitle=true,
  pdftitle={Constructive Repair in Realization Calculus},
  pdfauthor={Komorebi P.},
  pdfsubject={Obstructions, admissible corrections, and exact rational certificates}]{hyperref}
\usepackage{bookmark}
\usepackage[nameinlink,noabbrev]{cleveref}
\color{Ink}
\setlength{\parindent}{0pt}
\setlength{\parskip}{0.45em plus 0.1em minus 0.1em}
\linespread{1.075}
\setcounter{tocdepth}{1}
\setcounter{secnumdepth}{1}
\renewcommand{\bibname}{References}
\hyphenpenalty=250
\widowpenalty=10000
\clubpenalty=10000
\emergencystretch=2em
\raggedbottom
\setlist[itemize]{leftmargin=1.4em,itemsep=2pt,topsep=4pt}
\setlist[enumerate]{leftmargin=1.7em,itemsep=3pt,topsep=4pt}
\titleformat{\chapter}[display]
  {\normalfont\color{Pine}}{\sffamily\small\MakeUppercase{\chaptertitlename}\ \thechapter}
  {9pt}{\Huge\bfseries}
\titlespacing*{\chapter}{0pt}{12pt}{23pt}
\titleformat{\section}{\Large\bfseries\color{Pine}}{\thesection}{0.75em}{}
\titlespacing*{\section}{0pt}{17pt}{8pt}
\titleformat{\subsection}{\large\bfseries\color{Pine}}{}{0pt}{}
\titlespacing*{\subsection}{0pt}{13pt}{6pt}
\pagestyle{fancy}
\fancyhf{}
\fancyhead[L]{\sffamily\fontsize{8}{10}\selectfont\color{Slate}OPERATION KOMOREBI}
\fancyfoot[L]{\sffamily\fontsize{8}{10}\selectfont\color{Slate}REALIZATION CALCULUS}
\fancyfoot[R]{\sffamily\small\color{Pine}\thepage}
\renewcommand{\headrulewidth}{0.3pt}
\renewcommand{\footrulewidth}{0pt}
\fancypagestyle{plain}{\fancyhf{}\fancyfoot[L]{\sffamily\fontsize{8}{10}\selectfont\color{Slate}REALIZATION CALCULUS}\fancyfoot[R]{\sffamily\small\color{Pine}\thepage}\renewcommand{\headrulewidth}{0pt}}
\newtheoremstyle{quiet}{10pt}{7pt}{\normalfont}{}{\bfseries\color{Pine}}{.}{0.5em}{}
\theoremstyle{quiet}
\newtheorem{theorem}{Theorem}[chapter]
\newtheorem{proposition}[theorem]{Proposition}
\newtheorem{lemma}[theorem]{Lemma}
\newtheorem{corollary}[theorem]{Corollary}
\newtheorem{definition}[theorem]{Definition}
\newtheorem{example}[theorem]{Example}
\newtheorem{remark}[theorem]{Remark}
\renewcommand{\qedsymbol}{\ensuremath{\square}}
\newtcolorbox{scopebox}[1][]{enhanced,breakable,colback=PalePine,
  colframe=RuleGray,boxrule=0.3pt,arc=0pt,left=11pt,right=11pt,top=9pt,bottom=9pt,
  borderline west={1.5pt}{0pt}{Pine},fonttitle=\sffamily\bfseries,
  coltitle=Pine,colbacktitle=PalePine,attach title to upper={\par\smallskip},#1}
\newcommand{\kk}{\Bbbk}
\newcommand{\Q}{\mathbb{Q}}
\newcommand{\R}{\mathbb{R}}
\newcommand{\Z}{\mathbb{Z}}
\newcommand{\N}{\mathbb{N}}
\newcommand{\F}{\mathbb{F}}
\DeclareMathOperator{\im}{im}
\DeclareMathOperator{\coker}{coker}
\DeclareMathOperator{\Span}{span}
\DeclareMathOperator{\rank}{rank}
\DeclareMathOperator{\cost}{cost}
\DeclareMathOperator{\id}{id}
\DeclareMathOperator{\argmin}{argmin}
\newcommand{\Ob}{\mathcal{O}}
\newcommand{\code}[1]{\texttt{\detokenize{#1}}}
\newcommand{\decl}[1]{{\small\nolinkurl{#1}}}
\newcommand{\formal}[1]{\par\nobreak\begingroup\footnotesize\color{Slate}\sffamily Formal counterpart: \normalfont\nolinkurl{#1}.\par\endgroup}
\newcommand{\sourcescope}[1]{\par\begingroup\small\color{Slate}#1\par\endgroup}
\newcommand{\record}[3]{\noindent\begin{minipage}{\linewidth}\small\textbf{\nolinkurl{#1}}\hfill\textsf{line #2}\par\smallskip #3\end{minipage}\par\vspace{7pt}}
\newcommand{\moduleheading}[2]{\section{#1}\label{#2}\noindent\small Source: \nolinkurl{lean/KomorebiRepair/#1.lean}.\normalsize\par}
\begin{document}
\begin{titlepage}
\pagecolor{WarmPaper}
\color{Ink}
{\sffamily\fontsize{11}{15}\selectfont\color{Pine}OPERATION KOMOREBI\par}
\vspace{5pt}
{\sffamily\fontsize{9}{12}\selectfont\color{Slate}REALIZATION CALCULUS\par}
\vspace{1.1in}
{\sffamily\small\color{Pine}MATHEMATICAL MONOGRAPH\par}
\vspace{17pt}
{\fontsize{37}{42}\selectfont\bfseries Constructive Repair\par}
\vspace{7pt}
{\fontsize{27}{33}\selectfont in Realization Calculus\par}
\vspace{22pt}
{\color{Ochre}\rule{1.7in}{0.8pt}\par}
\vspace{19pt}
{\fontsize{16}{22}\selectfont Obstructions, admissible corrections,\par and exact rational certificates\par}
\vspace{0.55in}
\begin{minipage}{0.86\textwidth}
\large A correction is legitimate only when it removes the defect in the original mathematical object. Quotients identify what is missing. Admissibility determines which repairs are allowed.
\end{minipage}
\vfill
{\Large Komorebi P.\par}
\vspace{17pt}
{\sffamily\small\color{Slate}Release candidate 1.0.0rc3\quad /\quad 19 October 2026\par}
\vspace{6pt}
{\sffamily\footnotesize\color{Slate}51 unchanged Lean modules / fresh release audit\par}
\end{titlepage}
\nopagecolor
\setcounter{page}{2}
\chapter*{Orientation}
\addcontentsline{toc}{chapter}{Orientation}
This monograph develops a restricted theory of constructive repair. An original correction map is fixed, a family of defects is supplied, and additional correction directions may be introduced. The central question is not whether one can erase the description of a defect. It is whether the permitted changes produce an exact solution of the original correction equations.

The answer has three layers. Quotient obstructions decide feasibility. Linear sections attain the smallest possible added dimension over a field. A separate finite search chooses a minimum cost subset from a declared menu. The original finite-domain solver is retained, while the rational certificate extension checks exact coefficients without enumerating an infinite coefficient space. Faithful transport and conservation laws specify when a proposed simplification retains the information needed for those conclusions.

\begin{scopebox}[title=The central result]
Let $D:U\to V$ be linear over a field and let
\[
q:V\to V/\im D,\qquad O=\Span\{q(r_i):i\in I\}.
\]
If $V$ is finite dimensional, the least added dimension that repairs every $r_i$ is $\dim O$. If an allowed subspace $A\leq V$ satisfies $O\subseteq q(A)$, the minimum can be attained with all added directions in $A$. Selecting entries from a fixed menu is a different optimization problem.
\end{scopebox}

The mathematical development is self contained at the level of basic module theory, finite dimensional linear algebra, elementary arithmetic, and differentiation of polynomials. It refines the correction and obstruction viewpoint of Paper III rather than replacing the foundations of Realization Calculus. Its relation to the ambient realization program of Komorebi Space is discussed in \cref{sec:realization}.

\textbf{Formal scope.} This edition retains the original 49 theorems in ten modules and the later rational certificate and costed-menu layers. The structural tree certificates and branch producer are retained unchanged. The final appendix consolidates the installed interface, current evidence bindings, compact source provenance, and precisely limited support statement. Fresh audits remain separate from historical receipts.

\textbf{Reading route.} Chapters 1--3 cover correction and minimum dimension, Chapters 4--5 preservation and finite search, and Chapters 6--7 arithmetic and nonlinear boundaries. Chapters 8--10 develop rational certificates, explicit attainment, and costed menus. Chapter 11 gives the structural tree checker. The appendices record formal correspondence and verification. The final appendix describes the current release contract without duplicating chronological engineering reports. Declaration counts describe coverage, not novelty.

\clearpage
\tableofcontents

\chapter{Correction and quotient obstructions}\label{ch:basic}
\section{The original equation}
A residual describes a failure relative to a stated equation. A correction map describes the changes that are allowed to address that failure. These are separate pieces of data. Their separation prevents the phrase ``the obstruction'' from concealing an unstated choice of permissible operations.

Let $\kk$ be a commutative ring. All modules are unital $\kk$-modules and all maps in this chapter are $\kk$-linear. Fix
\[
D:U\longrightarrow V,
\]
and let $r:I\to V$ be a family of residuals. The index set $I$ need not be finite for the algebraic statements. When finite enumeration is used later, finiteness will be imposed explicitly.

\begin{definition}[Correction and repair]\label{def:repair}
A residual $r\in V$ is \emph{correctable} by $D$ when there is $u\in U$ such that
\begin{equation}\label{eq:correct}
Du=-r.
\end{equation}
An added correction map $E:W\to V$ \emph{repairs} $r$ when there are $u\in U$ and $w\in W$ with
\begin{equation}\label{eq:repair}
Du+Ew=-r.
\end{equation}
It repairs the family when this condition holds for every $r_i$. The witnesses may depend on $i$.
\end{definition}

The sign convention makes the repaired residual $r+Du+Ew$ equal to zero. In particular, the target of the equation remains $V$. A calculation in another space can establish \eqref{eq:repair} only if a valid route back to $V$ is supplied.

The family condition does not ask for one common pair $(u,w)$ that corrects every residual at once. It asks for one common added map $E$ and separate coefficients for each residual. Nor does it require those witnesses to vary continuously, measurably, or linearly with the index. Such conditions would define different problems.

\section{What the quotient remembers}
Write
\[
Q_D=\coker D=V/\im D,\qquad q_D:V\longrightarrow Q_D.
\]
The class $q_D(r)$ records the part of $r$ that cannot be changed by the original correction map. Since $q_D(Du)=0$, all residuals that differ by an original correction have the same class.

\begin{theorem}[Exact obstruction criterion]\label{thm:correctable}
For every $r\in V$,
\[
\exists u\in U\ (Du=-r)\quad\Longleftrightarrow\quad q_D(r)=0.
\]
\end{theorem}
\begin{proof}
Applying $q_D$ to $Du=-r$ gives $0=-q_D(r)$, hence $q_D(r)=0$. Conversely, $q_D(r)=0$ means $r\in\im D$. Choose $v\in U$ with $Dv=r$. Then $D(-v)=-r$.
\end{proof}
\formal{correctable_iff_obstruction_zero}

This is an exact criterion for the declared linear correction equation. It is not a universal test for a nonlinear equation from which $D$ might have been obtained by differentiation. Chapter 7 gives counterexamples in both directions.

The quotient also has a universal property. If $L:V\to Z$ annihilates $\im D$, it descends to a unique linear map $\overline L:Q_D\to Z$ with $L=\overline Lq_D$. Thus every linear detector that ignores correctable directions factors through the same quotient. This is the classical cokernel construction used in Paper III~\cite{paperiii}. Canonicality is relative to $D$, not to an unspecified mathematical problem.

\begin{example}[A visible residual]\label{ex:plane}
Let $D:\Q\to\Q^2$ be $D(u)=(u,0)$ and let $r=(2,3)$. The quotient is identified with the second coordinate. The correction $u=-2$ removes the first coordinate but leaves $(0,3)$. The functional $L(x,y)=y$ satisfies $LD=0$ and $L(r)=3$, so no original correction eliminates $r$.
\end{example}

\section{Repair as an image condition}
An added map $E:W\to V$ contributes precisely the classes in $\im(q_DE)$. This leads to the basic feasibility criterion.

\begin{theorem}[A single residual]\label{thm:single}
For every $E:W\to V$ and $r\in V$,
\begin{equation}\label{eq:single}
\exists u,w\ (Du+Ew=-r)
\quad\Longleftrightarrow\quad
q_D(r)\in\im(q_DE).
\end{equation}
\end{theorem}
\begin{proof}
Suppose $Du+Ew=-r$. Applying the quotient gives $q_D(Ew)=-q_D(r)$, so $q_D(r)=q_D(E(-w))$ lies in the image.

Conversely, choose $w$ with $q_D(Ew)=q_D(r)$. Then $q_D(r-Ew)=0$, so there is $v\in U$ with $Dv=r-Ew$. Taking $u=-v$ and $w'=-w$ gives
\[
Du+Ew'=-(r-Ew)-Ew=-r.
\]
The witnesses solve the equation in $V$, not just in its quotient.
\end{proof}
\formal{repairable_iff_obstruction_mem_range}

\Needspace{8\baselineskip}
\begin{corollary}[Reconstruction from a quotient witness]\label{cor:reconstruct}
If $q_D(Ew)=q_D(r)$, there are $u\in U$ and $w'\in W$ such that
\[
Du+Ew'+r=0.
\]
\end{corollary}
\begin{proof}
Use the witnesses constructed in the reverse direction of \cref{thm:single} and move $r$ to the left.
\end{proof}
\formal{reconstruct_cancellation}

\section{The obstruction span of a family}
\begin{definition}[Required obstruction span]\label{def:span}
For a residual family $r:I\to V$, put
\[
O_D(r)=\Span_{\kk}\{q_D(r_i):i\in I\}\leq Q_D.
\]
We write $O$ when the family and correction map are understood.
\end{definition}

Every $q_D(r_i)$ lies in $O$. The use of a span matters because an added linear correction map supplies a submodule of quotient classes, not merely a list of unrelated vectors.

\begin{theorem}[Family feasibility]\label{thm:family}
The map $E:W\to V$ repairs the family $r$ if and only if
\begin{equation}\label{eq:family}
O_D(r)\subseteq\im(q_DE).
\end{equation}
\end{theorem}
\begin{proof}
If every residual is repairable, \cref{thm:single} places every generator $q_D(r_i)$ in $\im(q_DE)$. The image is a submodule, so it contains their span. Conversely, \eqref{eq:family} places each generator in that image, and \cref{thm:single} supplies its correction witnesses.
\end{proof}
\formal{family_repairable_iff_span_le}

For an empty family, $O=0$ and the family condition is vacuous. For a nonempty family whose residuals are already correctable, $O$ is still zero. Both cases should permit a repair of added dimension zero. The next chapter includes them without a special exception.

\begin{scopebox}[title=What has been established]
The quotient determines feasibility only because the correction map and the original equation have been declared. A positive quotient witness can be lifted to actual cancellation in $V$. A disappearance caused by forgetting part of $V$ is a different operation and requires a preservation theorem.
\end{scopebox}

\chapter{Minimum repair dimension}\label{ch:minimal}
\section{A lower bound that every repair must satisfy}
In this chapter $\kk$ is a field. The added space $W$ is finite dimensional when its numerical dimension is compared. The attainment result itself uses a linear section and does not require all ambient spaces to be finite dimensional.

\begin{theorem}[Dimension lower bound]\label{thm:lower}
If $E:W\to V$ repairs $r$ and $W$ is finite dimensional, then $O=O_D(r)$ is finite dimensional and
\[
\dim_{\kk}O\leq\dim_{\kk}W.
\]
\end{theorem}
\begin{proof}
By \cref{thm:family}, $O\subseteq\im(q_DE)$. The image of a linear map from a finite dimensional space is finite dimensional and has dimension at most that of its domain. Hence
\[
\dim O\leq\dim\im(q_DE)\leq\dim W.\qedhere
\]
\end{proof}
\formal{repair_dimension_lower_bound}

The bound measures missing independent directions, not the number of residuals. A large family can have a one dimensional obstruction span. Conversely, a single residual never requires more than one freely chosen direction, since its own negative may be made available. These statements concern unconstrained linear corrections and say nothing about the cost of producing a permitted direction in a domain application.

\section{An attaining construction}
\begin{lemma}[A section of the quotient]\label{lem:section}
There is a linear map $s:Q_D\to V$ with
\[
q_Ds=\id_{Q_D}.
\]
\end{lemma}
\begin{proof}
Choose a basis of the vector space $Q_D$. Since $q_D$ is surjective, choose a representative in $V$ of each basis vector. Extend this assignment linearly. Its composite with $q_D$ fixes the basis and therefore fixes all of $Q_D$.
\end{proof}
\formal{exists_obstruction_section}

This choice of section is generally nonunique. Changing a representative by a vector in $\im D$ produces another valid section. The construction is an existence proof using classical choice, not an algorithm for selecting a basis.

\begin{theorem}[Attainment on the obstruction span]\label{thm:attain}
Let $\iota_O:O\hookrightarrow Q_D$ be inclusion. There exists a linear map $E:O\to V$ such that
\begin{equation}\label{eq:section-on-O}
q_DE=\iota_O
\end{equation}
and $E$ repairs the family.
\end{theorem}
\begin{proof}
Choose the section $s$ from \cref{lem:section} and set $E=s\iota_O$. Then $q_DE=q_Ds\iota_O=\iota_O$. Every $z\in O$ is therefore in $\im(q_DE)$. Family feasibility follows from \cref{thm:family}.
\end{proof}
\formal{exists_attaining_repair}

\[
\begin{tikzcd}[column sep=large,row sep=large]
O \arrow[r,"E"] \arrow[dr,hook,"\iota_O"'] & V \arrow[d,"q_D"]\\
& Q_D
\end{tikzcd}
\]
The triangle records exactly what is preserved: a required quotient class is lifted to a representative of that same class. The lift does not manufacture an unrelated object with a similar invariant.

\begin{theorem}[The attained free minimum]\label{thm:minimum}
If $V$ is finite dimensional, the minimum possible dimension of a finite dimensional added space that repairs the family is
\begin{equation}\label{eq:minimum}
\min\bigl\{\dim W: E:W\to V\text{ repairs }r\bigr\}=\dim O_D(r).
\end{equation}
\end{theorem}
\begin{proof}
Every admissible term on the left is at least $\dim O$ by \cref{thm:lower}. Because $V$ is finite dimensional, its quotient and the subspace $O$ are finite dimensional. The map in \cref{thm:attain} is a repair with domain $O$, so the bound is attained.
\end{proof}
\formal{free_minimal_repair_rank}

No desired existence statement occurs among the inputs. The theorem proves nonemptiness of the optimization class as well as its lower bound.

\section{Directions restricted to a subspace}
An application usually prohibits some corrections. A first useful constraint is a subspace $A\leq V$ containing every permitted added direction. Its quotient image is
\[
q_D(A)=\im\bigl(q_D\vert_A:A\to Q_D\bigr).
\]

\begin{theorem}[Attainment inside a legal subspace]\label{thm:legal}
If $O\subseteq q_D(A)$, there exists a linear map $E:O\to V$ satisfying
\[
\im E\subseteq A,\qquad q_DE=\iota_O,
\]
and repairing the entire family. If $V$ is finite dimensional, its added dimension is optimal even among all unrestricted repairs.
\end{theorem}
\begin{proof}
Consider the surjection
\[
f:A\longrightarrow q_D(A),\qquad a\longmapsto q_D(a).
\]
Choose a linear section $s_A:q_D(A)\to A$. Let $j:O\hookrightarrow q_D(A)$ and $\iota_A:A\hookrightarrow V$ be inclusions. Define
\[
E=\iota_As_Aj.
\]
Its image is in $A$ by construction. Moreover, $q_D\iota_As_Aj=j$ as a map into $Q_D$, which is precisely $\iota_O$. By \cref{thm:family}, $E$ repairs the family. The lower bound of \cref{thm:lower} proves optimality when dimensions are finite.
\end{proof}
\formal{admissible_subspace_attainment}

\begin{theorem}[Subspace impossibility]\label{thm:legal-no-go}
If $O\not\subseteq q_D(A)$, no added map $E:W\to V$ with image contained in $A$ repairs the family.
\end{theorem}
\begin{proof}
Such a map satisfies $\im(q_DE)\subseteq q_D(A)$. If it repaired the family, \cref{thm:family} would imply $O\subseteq\im(q_DE)\subseteq q_D(A)$, contrary to the hypothesis.
\end{proof}
\formal{admissible_subspace_no_go}

The word ``legal'' means membership in a specified linear subspace. It does not mean that the paper has constructed an arithmetic host, an equivariant representation, or a topology on that space. Those interpretations require separate witnesses.

\section{Matrix coordinates and a boundary over rings}
For finite matrices, the abstract minimum has a familiar coordinate form.

\begin{corollary}[Rank difference in matrix coordinates]\label{cor:rank-difference}
Let the columns of a matrix $\mathsf R$ be a finite residual family and let $\mathsf D$ represent $D$. Then
\[
\dim O_D(r)=\rank[\mathsf D\mid\mathsf R]-\rank\mathsf D.
\]
\end{corollary}
\begin{proof}
Set $M=\im D+\Span\{r_i\}$. The restriction $q_D\vert_M$ has image $O$ and kernel $\im D$. Rank and nullity give $\dim M=\dim\im D+\dim O$. The columns of $[\mathsf D\mid\mathsf R]$ span $M$.
\end{proof}
\sourcescope{This coordinate corollary is an ordinary derivation included for the computational interpretation. The v0.2.0 Lean release proves the abstract attained dimension theorem, but does not yet identify this matrix formula with the Python rank implementation.}

The field hypothesis in the section argument is substantive. Over $\Z$, take $D:\Z\to\Z$ to be multiplication by two. The quotient map $q:\Z\to\Z/2\Z$ has no $\Z$-linear section. A section would send the nonzero class to an integer $n$ with $2n=0$, forcing $n=0$ and contradicting $q(n)=1$. This elementary example explains why the module feasibility theorem does not automatically imply a splitting theorem over arbitrary rings.



\chapter{Finite menus and conservation}\label{ch:menu}
\section{A menu is an actual coefficient map}
Return to a commutative ring $\kk$ and modules $U,V$. Let $e:J\to V$ be a family of candidate directions and let $S\subseteq J$ be finite. A coefficient vector is a function $a:S\to\kk$. Define
\begin{equation}\label{eq:menu}
E_S:\kk^S\longrightarrow V,\qquad E_S(a)=\sum_{j\in S}a_je_j.
\end{equation}
Finite addition and scalar multiplication show that $E_S$ is linear. The coefficient domain is the selected set itself, so no unselected direction appears in the equation.

Write
\[
O_S=\Span_{\kk}\{q_D(e_j):j\in S\}.
\]

\begin{lemma}[Quotient image of a menu]\label{lem:menu-range}
One has $\im(q_DE_S)=O_S$.
\end{lemma}
\begin{proof}
For every coefficient vector $a$,
\[
q_D(E_Sa)=\sum_{j\in S}a_jq_D(e_j),
\]
which lies in $O_S$. Conversely, every generator $q_D(e_j)$ is the image of the coefficient vector that equals one at $j$ and zero elsewhere. Since the image is a submodule, it contains their span.
\end{proof}
\formal{menu_quotient_range}

\begin{theorem}[Exact finite menu feasibility]\label{thm:menu}
The selected menu $S$ repairs the family exactly when
\begin{equation}\label{eq:menu-criterion}
O_D(r)\subseteq O_S.
\end{equation}
Equivalently,
\[
\forall i\in I\ \exists u_i\in U\ \exists a_i\in\kk^S,
\qquad Du_i+\sum_{j\in S}a_{ij}e_j=-r_i.
\]
\end{theorem}
\begin{proof}
Apply \cref{thm:family} to the actual map $E_S$ and use \cref{lem:menu-range}. Its reverse direction supplies the original correction witnesses and the selected coefficients, rather than only an abstract statement about dimensions.
\end{proof}
\formal{finite_menu_repair_iff_span}

\Needspace{8\baselineskip}
\begin{proposition}[Monotonicity of feasibility]\label{prop:menu-mono}
If $S\subseteq T$, then $O_S\subseteq O_T$. Every family repairable with $S$ is therefore repairable with $T$.
\end{proposition}
\begin{proof}
Each generator defining $O_S$ is also a generator defining $O_T$. Apply \cref{thm:menu}. Equivalently, extend each coefficient vector by zero on $T\setminus S$.
\end{proof}
\formal{repairable_menu_mono}

An external admissibility condition need not be monotone. A rule allowing at most one selected direction can reject an enlarged menu even though algebraic feasibility only improves. Cost and admissibility must not be confused with span inclusion.

\section{One obstruction, two selected entries}
\begin{example}[The difference between the two minima]\label{ex:menu}
Over $\Q$, let
\[
D(u)=(u,0,0),\qquad r=(0,1,0),
\]
and make only the following directions available:
\[
e_0=(0,1,1),\qquad e_1=(0,0,1).
\]
The quotient by $\im D$ is the plane of the second and third coordinates. There $q_D(r)=(1,0)$, $q_D(e_0)=(1,1)$, and $q_D(e_1)=(0,1)$.
\end{example}

Using only $e_0$, the repair equation reads
\[
(u,a,a)=(0,-1,0).
\]
The second coordinate forces $a=-1$ and the third forces $a=0$, a contradiction. Using only $e_1$, the second coordinate stays zero and cannot become $-1$. Both directions together give the exact witness
\begin{equation}\label{eq:two-entry-witness}
D(0)-e_0+e_1=(0,-1,0)=-r.
\end{equation}
The empty menu also fails. Its correction range is the first coordinate axis.

\begin{center}
\begin{tabular}{@{}lll@{}}
\toprule
Selected set & Quotient directions & Repairs $r$?\\
\midrule
$\varnothing$ & $0$ & No\\
$\{0\}$ & $\Span\{(1,1)\}$ & No\\
$\{1\}$ & $\Span\{(0,1)\}$ & No\\
$\{0,1\}$ & $\Q^2$ & Yes\\
\bottomrule
\end{tabular}
\end{center}

The obstruction span has dimension one. A freely chosen direction $r$ repairs the residual with coefficient $-1$. The prescribed menu nevertheless requires two selected entries. Both statements are formalized in the examples module. The difference is not a contradiction: one cost counts independent added dimensions, the other counts named entries that must be purchased or enabled.

A weighting makes the distinction sharper. Assigning positive costs $c_0,c_1$ gives optimum $c_0+c_1$ for this menu. A subspace containing $e_0-e_1$ does not charge those two entries separately unless its cost model explicitly does so.

\section{Preserved linear invariants}
Conservation laws supply impossibility certificates without enumerating every correction.

\begin{proposition}[Invariant preservation]\label{prop:invariant}
Let $L:V\to Z$ be linear. If $LD=0$ and $LE=0$, then
\[
L(v+Du+Ew)=L(v)
\]
for all $v,u,w$.
\end{proposition}
\begin{proof}
Linearity gives $L(v+Du+Ew)=L(v)+LD(u)+LE(w)$. The last two terms vanish.
\end{proof}
\formal{invariant_preserved}

\begin{theorem}[Conservation obstruction]\label{thm:conservation}
If $LD=0$, $LE=0$, and $L(r)\ne0$, then $E$ does not repair $r$.
\end{theorem}
\begin{proof}
A repaired residual would be $r+Du+Ew=0$. Applying \cref{prop:invariant} would give $L(r)=L(0)=0$, contrary to the hypothesis.
\end{proof}
\formal{conservation_no_go}

\begin{corollary}[Conservation for a selected menu]\label{cor:menu-conservation}
Suppose $LD=0$ and $L(e_j)=0$ for every $j\in S$. If $L(r)\ne0$, no coefficients from $S$ repair $r$.
\end{corollary}
\begin{proof}
For every $a:S\to\kk$,
\[
L(E_Sa)=\sum_{j\in S}a_jL(e_j)=0.
\]
Thus $LE_S=0$, and \cref{thm:conservation} applies.
\end{proof}
\formal{conservation_menu_no_go}

For instance, let $D(u)=(u,0,0)$ and permit additional directions only on the third coordinate axis. The invariant $L(x,y,z)=y$ annihilates every legal correction. It detects $r=(0,1,0)$, so no legal repair exists. Enlarging the number of copies of the third coordinate direction cannot help.

The theorem is only as meaningful as the invariant it uses. In a domain application one must establish that $L$ is the intended native invariant and that every permitted transformation really preserves it. An abstract functional with a suggestive name is not such a proof.

\chapter{Faithful transport of obstructions}\label{ch:transport}
\section{The induced quotient map}
Suppose two correction systems are related by a commuting square
\begin{equation}\label{eq:square}
\begin{tikzcd}[column sep=large,row sep=large]
U \arrow[r,"D"] \arrow[d,"a"'] & V \arrow[d,"b"]\\
U' \arrow[r,"D'"'] & V'.
\end{tikzcd}
\qquad bD=D'a.
\end{equation}
All four spaces are modules over the same commutative ring. Denote the quotient maps by $q_D$ and $q_{D'}$.

\begin{lemma}[The correction image descends]\label{lem:range-maps}
The map $b$ sends $\im D$ into $\im D'$.
\end{lemma}
\begin{proof}
If $v=Du$, then $b(v)=bD(u)=D'a(u)$, which lies in $\im D'$.
\end{proof}
\formal{correction_range_maps}

It follows that there is a well defined linear map
\[
B:Q_D\longrightarrow Q_{D'},\qquad B(q_D(v))=q_{D'}(b(v)).
\]
Indeed, changing $v$ by an element of $\im D$ changes $b(v)$ by an element of $\im D'$. This is the quotient map constructed in the formal source, not an assumed function on unidentified obstruction spaces.

\begin{proposition}[Naturality]\label{prop:naturality}
For every residual $r\in V$,
\begin{equation}\label{eq:naturality}
B(q_D(r))=q_{D'}(b(r)).
\end{equation}
\end{proposition}
\begin{proof}
This is the defining equation of the induced quotient map.
\end{proof}
\formal{transport_obstruction_naturality}

Naturality transports a correctable residual to a correctable one. The converse can fail because $B$ may annihilate a nonzero class. A second condition is needed to use target vanishing as a source certificate.

\section{Faithfulness on precisely the required data}
\begin{definition}[Faithfulness on a submodule]\label{def:faithful}
For $O\leq Q_D$, the map $B$ is \emph{faithful on $O$} when
\[
\forall z\in O,\qquad Bz=0\Longrightarrow z=0.
\]
\end{definition}

This condition is equivalent to $\ker B\cap O=0$. It does not require injectivity on all of $Q_D$ when only a smaller family of obstruction classes matters.

\begin{proposition}[Equivalent injectivity criterion]\label{prop:faithful-inj}
The map $B$ is faithful on $O$ if and only if its restriction $B\vert_O:O\to Q_{D'}$ is injective.
\end{proposition}
\begin{proof}
If $B$ is faithful and $Bx=By$ for $x,y\in O$, then $x-y\in O$ and $B(x-y)=0$. Faithfulness gives $x=y$. Conversely, if $B\vert_O$ is injective and $Bz=0=B0$ for $z\in O$, injectivity gives $z=0$.
\end{proof}
\formal{faithfulOn_iff_injective}

\begin{theorem}[Reflection of correctability]\label{thm:faithful}
Let $O=O_D(r)$. If $B$ is faithful on $O$, then for every $i\in I$,
\[
b(r_i)\text{ is correctable by }D'
\quad\Longleftrightarrow\quad
r_i\text{ is correctable by }D.
\]
\end{theorem}
\begin{proof}
By \cref{thm:correctable}, target correctability means $q_{D'}(b(r_i))=0$. Naturality identifies this class with $B(q_D(r_i))$. Because $q_D(r_i)\in O$, faithfulness then yields $q_D(r_i)=0$, which is source correctability. Conversely, if $q_D(r_i)=0$, naturality and linearity give $q_{D'}(b(r_i))=B0=0$.
\end{proof}
\formal{transport_faithfulness_on_obstruction_span}

\begin{proposition}[A witness to information loss]\label{prop:loss-witness}
If $B$ is not faithful on $O$, there is $z\in O$ such that
\[
z\ne0,\qquad Bz=0.
\]
\end{proposition}
\begin{proof}
Negating the quantified implication in \cref{def:faithful} gives an element of $O$ whose image is zero and which is itself nonzero.
\end{proof}
\formal{nonfaithful_witness}

The witness may be a linear combination of the supplied obstruction classes rather than one particular generator. This is why injectivity on the entire required span is a stronger preservation test than checking that a short list of nonzero classes individually remains nonzero.

\section{A transport that conceals failure}
\begin{example}[Correct in the target, false in the source]\label{ex:loss}
Let
\[
D:\Q\to\Q^2,\quad D(u)=(u,0),\qquad
D'=\id_{\Q},\quad a=\id_{\Q},\quad b(x,y)=x.
\]
Then $bD=D'a$. For $r=(0,1)$ there is no $u$ with $(u,0)=(0,-1)$, so the source residual is not correctable. Its image $b(r)=0$ is correctable by $D'$.
\end{example}

Here $Q_D$ is one dimensional and $Q_{D'}=0$. The induced map $B$ annihilates every source obstruction. It is therefore not faithful on the span of $q_D(r)$. The formal example proves both the false positive produced by projection and the failure of the faithfulness hypothesis.

A second elementary calculation clarifies why faithfulness is imposed on a span. The map $B:\Q^2\to\Q$, $B(x,y)=x+y$, sends each standard basis vector to the nonzero number one. Nevertheless, it kills their difference $(1,-1)$. Survival of generators is not injectivity on their span. This calculation is explanatory and has no separate declaration in the present formal source.

\begin{scopebox}[title=Transport is not repair]
A repair changes the allowed correction image while keeping the original equation meaningful. A transport changes the representation of that equation. The two may be composed, but vanishing in the transported space certifies source correctability only when the stated preservation hypotheses hold.
\end{scopebox}

\chapter{Certified finite search}\label{ch:search}
\section{The finite contract}
The algebraic existence theorems need not be executable. A linear section of a quotient may be chosen by classical choice, and a feasibility statement over an arbitrary infinite field need not come with a terminating decision procedure. The finite reference solver uses a different contract.

Assume $\kk$, $U$, $I$, and $J$ have specified finite enumerations. The ring $\kk$ is commutative, $U,V$ are $\kk$-modules, and equality in $V$ is decidable. The target module $V$ need not itself have a finite enumeration. Fix a linear order on $J$ so that the candidate universe can be enumerated in a definite order.

For each finite subset $S\subseteq J$, define the Boolean test
\begin{equation}\label{eq:boolean-check}
\operatorname{check}(S)=\operatorname{decide}\left(
\forall i\in I\ \exists u\in U\ \exists a\in\kk^S,\quad Du+E_Sa=-r_i
\right).
\end{equation}
All quantified domains are finite. Decidable equality in $V$ determines each candidate equation.

\begin{proposition}[The Boolean test has the intended meaning]\label{prop:check}
The equality $\operatorname{check}(S)=\mathrm{true}$ holds if and only if $E_S$ repairs the residual family.
\end{proposition}
\begin{proof}
The Boolean is defined by deciding precisely the quantified equations in \cref{def:repair} for the map $E_S$. No intermediate rank guess or unproved feasibility oracle replaces those equations.
\end{proof}
\formal{repairMenuCheck_spec}

Additional admissibility is a supplied function $\operatorname{allowed}:\mathcal P(J)\to\{\mathrm{false},\mathrm{true}\}$. The cost is any supplied function $c:\mathcal P(J)\to\N$. The formal theorem does not require $c$ to be additive in the selected entries.

\section{Enumerating menus and selecting an optimum}
Enumerate $J$ as a sorted list. Take every sublist and convert it to a finite set. The resulting list, denoted $\mathcal L_J$, contains every subset of $J$.

\begin{lemma}[Coverage of the menu universe]\label{lem:allmenus}
Every finite subset $S\subseteq J$ occurs in $\mathcal L_J$.
\end{lemma}
\begin{proof}
Filter the sorted full list by membership in $S$. The result is a sublist of the full list and therefore appears among its sublists. Converting this filtered list to a finite set gives exactly $S$.
\end{proof}
\formal{mem_allMenus}

For a finite list $\mathcal L$ of candidates, a Boolean predicate $h$, and a natural cost $c$, define $\operatorname{bestByCheck}(h,c,\mathcal L)$ to filter by $h$ and then return an argument of minimum cost. It returns \code{none} exactly when the filtered list is empty. Ties need not define a mathematically unique optimum.

\begin{lemma}[Finite selection]\label{lem:best}
If this selector returns $x$, then $x\in\mathcal L$, $h(x)=\mathrm{true}$, and $c(x)\leq c(y)$ for every $y\in\mathcal L$ with $h(y)=\mathrm{true}$. It returns \code{none} if and only if no listed candidate passes $h$.
\end{lemma}
\begin{proof}
Every element of the filtered list is an original list member satisfying $h$. A minimum cost element of that finite list has cost no larger than that of any other member. The selector returns no element exactly when the filtered list is empty, which is exactly the stated rejection condition.
\end{proof}
\formal{bestByCheck_spec}

The repair optimizer is
\[
\operatorname{bestRepairMenu}
=\operatorname{bestByCheck}\bigl(
S\mapsto\operatorname{allowed}(S)\wedge\operatorname{check}(S),\ c,\ \mathcal L_J\bigr).
\]

\begin{theorem}[Soundness and minimum cost]\label{thm:optimizer}
If the optimizer returns \code{some S}, then $S$ is allowed, $E_S$ repairs the family, and
\[
c(S)\leq c(T)
\]
for every allowed subset $T\subseteq J$ whose menu map repairs the family.
\end{theorem}
\begin{proof}
Apply \cref{lem:best}. Passing the conjunction gives both admissibility and $\operatorname{check}(S)=\mathrm{true}$. Proposition~\ref{prop:check} converts the latter into actual repair feasibility. Every competing $T$ occurs in $\mathcal L_J$ by \cref{lem:allmenus}, so the same finite selection lemma compares its cost to $c(S)$.
\end{proof}
\formal{menu_optimality_certificate}

\begin{theorem}[Complete rejection within the declared universe]\label{thm:none}
The optimizer returns \code{none} if and only if no allowed subset of $J$ repairs the family.
\end{theorem}
\begin{proof}
If it returns \code{none}, \cref{lem:best} excludes every passing member of $\mathcal L_J$. Any allowed feasible subset would occur there by \cref{lem:allmenus} and pass the Boolean test by \cref{prop:check}, a contradiction. Conversely, if no allowed subset repairs the family, no listed subset satisfies the filter, so the selector returns \code{none}.
\end{proof}
\formal{menu_search_none_iff}

The output is a selected subset or a rejection, not an extracted table of correction coefficients. The correctness theorem establishes existence of the coefficients. A separately implemented witness extraction or serialized certificate checker would provide those coefficients as a returned data object.

\section{A fully evaluated example over two elements}\label{sec:f2}
Take the data of \cref{ex:menu} over $\F_2$. Negation is now the identity. The equation is
\[
(u,a,a+b)=(0,1,0).
\]
Thus the full menu succeeds with $u=0$, $a=b=1$. The two singleton menus still fail.

\begin{center}
\begin{tabular}{@{}lcll@{}}
\toprule
Menu & Cardinality & All menus allowed & At most one entry\\
\midrule
$\varnothing$ & 0 & Infeasible & Infeasible\\
$\{0\}$ & 1 & Infeasible & Infeasible\\
$\{1\}$ & 1 & Infeasible & Infeasible\\
$\{0,1\}$ & 2 & Feasible & Not allowed\\
\bottomrule
\end{tabular}
\end{center}

With cardinality as the cost, the two evaluations are
\begin{Verbatim}[fontsize=\small,frame=single,framesep=8pt,rulecolor=\color{RuleGray}]
All subsets permitted:                  some {0, 1}
Only subsets with cardinality at most 1: none
\end{Verbatim}
These equalities are the statements of
\decl{executable_menu_returns_both} and
\decl{executable_restricted_menu_rejects}. Their source proofs use kernel reduction. The generic optimizer theorem then proves that every feasible subset has cardinality at least two, recorded as \decl{executable_menu_is_globally_optimal}. Printed evaluation output is kept distinct from the proof of that output.

\section{What the algorithm does not promise}
The reference solver is deliberately finite and exhaustive. With $t=|J|$, there are $2^t$ candidate subsets. For a subset of size $s$, each residual has at most $|U|\,|\kk|^s$ candidate coefficient and correction assignments to examine. Summing this upper bound over all subsets gives
\[
|I|\,|U|\sum_{s=0}^t\binom ts|\kk|^s
=|I|\,|U|(1+|\kk|)^t
\]
possible equation tests before accounting for short circuiting and inadmissible subsets. This is an elementary enumeration bound, not a formal runtime theorem about the compiled Lean implementation.

The Python prototype uses exact Gaussian elimination for its linear equations. The finite Lean solver instead decides those equations by finite enumeration. They share a mathematical specification in the relevant finite cases, but the present release does not contain a proof identifying their implementations. In particular, the Lean optimizer is not an executable search over arbitrary rational coefficient assignments.

\chapter{Arithmetic digit repair}\label{ch:arithmetic}
\section{An exact correction at one scale}
Arithmetic repair adds a second kind of constraint: the new correction must improve divisibility while retaining the previous approximation. A quadratic equation makes the mechanism visible without appealing to a general existence theorem.

The results in this chapter are first stated over $\Z$. An integer $p$ need not be prime for the finite divisibility implications. The interpretation as a $p$-adic lift uses a positive prime $p$.

\begin{lemma}[A linear digit from B\'ezout data]\label{lem:digit}
Let $p,A,c\in\Z$. If there are $s,b\in\Z$ with $sA+bp=1$, then there is $t\in\Z$ such that
\[
p\mid c+At.
\]
\end{lemma}
\begin{proof}
Set $t=-sc$. Then
\[
c+At=c(1-As)=pbc.\qedhere
\]
\end{proof}
\formal{Arithmetic.linear_digit_exists}

\begin{lemma}[Quadratic correction identity]\label{lem:quad}
Let $p,m,x,a,c,t\in\Z$ satisfy
\[
x^2-a=mc,\qquad p\mid c+2xt,\qquad p\mid m.
\]
Then
\[
mp\mid (x+mt)^2-a.
\]
\end{lemma}
\begin{proof}
Write $c+2xt=pd$ and $m=ps$. Direct expansion gives
\begin{align*}
(x+mt)^2-a
&=(x^2-a)+2mxt+m^2t^2\\
&=m(c+2xt)+m^2t^2\\
&=mp(d+st^2).\qedhere
\end{align*}
The last expression displays the required quotient explicitly.
\end{proof}
\formal{Arithmetic.quadratic_digit_correction}

The term $m^2t^2$ is not discarded as an approximation. Its divisibility follows from $p\mid m$. This is what makes a linear digit equation sufficient for the finite quadratic step.

\section{The quadratic lifting theorem}
\begin{theorem}[One higher precision]\label{thm:hensel}
Let $p,x,a\in\Z$ and $k\in\N$ with $k\geq1$. Assume
\[
p^k\mid x^2-a,\qquad \gcd(2x,p)=1.
\]
Then there exists $t\in\Z$ such that
\begin{equation}\label{eq:hensel}
p^{k+1}\mid (x+p^kt)^2-a.
\end{equation}
\end{theorem}
\begin{proof}
Choose $c\in\Z$ with $x^2-a=p^kc$. Coprimality supplies B\'ezout data for $2x$ and $p$, so \cref{lem:digit} gives $t$ with $p\mid c+2xt$. Since $k\geq1$, one has $p\mid p^k$. Apply \cref{lem:quad} with $m=p^k$.
\end{proof}
\formal{Arithmetic.quadratic_hensel_step}

The formal hypothesis is \code{IsCoprime (2 * x) p}, expressed by B\'ezout data. The theorem constructs an integer representative of a correction class. It does not state that this representative is already the unique digit between zero and $p-1$, nor that an infinite succession of such choices has been constructed in a complete field.

\section{An exact seven-adic trace}\label{sec:trace}
For $x^2=2$ at $p=7$, start with $x_1=3$. At each step, let
\[
c_k=\frac{x_k^2-2}{7^k},\qquad
c_k+2x_kt_k\equiv0\pmod7,\qquad
x_{k+1}=x_k+7^kt_k.
\]
The following table records normalized choices for this example. Every displayed digit equation is an exact integer calculation.

\begin{center}
\begin{tabular}{@{}rrrrr@{}}
\toprule
$k$ & $x_k$ & $c_k\bmod7$ & $t_k$ & $x_{k+1}$\\
\midrule
1 & 3 & 1 & 1 & 10\\
2 & 10 & 2 & 2 & 108\\
3 & 108 & 6 & 6 & 2,166\\
4 & 2,166 & 1 & 1 & 4,567\\
5 & 4,567 & 2 & 2 & 38,181\\
6 & 38,181 & 1 & 1 & 155,830\\
\bottomrule
\end{tabular}
\end{center}

All representatives are congruent to three modulo seven, so $2x_k\equiv6\pmod7$ remains a unit. The formal statement \decl{Arithmetic.hensel_seven_trace} verifies the seven congruences
\[
7^k\mid x_k^2-2\qquad (1\leq k\leq7)
\]
for the displayed representatives. It is a finite conjunction of exact divisibility facts. The general one step theorem is proved separately.

For example,
\[
10^2-2=98=7^2\cdot2,
\]
and the next digit satisfies $2+2\cdot10\cdot2=42$. Hence
\[
108=10+7^2\cdot2,
\qquad 108^2-2=7^3\cdot34.
\]
The new digit corrects the residual quotient at precisely the next required modulus.

\section{Why a singular derivative is inconclusive}
\begin{proposition}[A modular solution that does not lift]\label{prop:mod4}
The equation $x^2=2$ has no solution in $\Z/4\Z$.
\end{proposition}
\begin{proof}
The four residue classes have squares $0,1,0,1$, respectively. None has square two.
\end{proof}
\formal{Arithmetic.singular_mod_four_impossible}

Nevertheless, $x=0$ solves the equation modulo two. Its two possible lifts modulo four are zero and two, and both fail. Thus existence at one precision is not itself an existence theorem at the next.

\begin{proposition}[Singularity does not prohibit lifting]\label{prop:singular-success}
For every $k\in\N$,
\[
2^k\mid 1^2-1.
\]
\end{proposition}
\begin{proof}
The residual is zero, which is divisible by every integer.
\end{proof}
\formal{Arithmetic.singular_derivative_can_still_lift}

For the equation $x^2-1=0$, the derivative at $x=1$ is two and is not a unit modulo two. The exact root still persists at every dyadic precision. A unit derivative is a sufficient hypothesis for \cref{thm:hensel}, not a necessary condition for every possible lift.

\chapter{Nonlinear limits and native realizations}\label{ch:limits}
\section{A linearized problem is a different equation}
Suppose a nonlinear map $F$ is differentiable at an approximation $x_0$. Its linearized correction equation is
\begin{equation}\label{eq:linearization}
DF(x_0)u=-F(x_0).
\end{equation}
The cokernel criterion completely decides \eqref{eq:linearization} once the derivative is fixed. It does not, by itself, decide whether there is $x$ with $F(x)=0$. Higher order terms are absent from the linearized equation, and they may either create a solution or prevent the linear prediction from reaching one.

The next two propositions are not warnings based only on examples computed numerically. The formal source includes actual derivative statements and exact proofs of the contrary conclusions.

\begin{proposition}[A nonlinear root with a nonzero first order obstruction]\label{prop:nonlinear-root}
For $a\in\R\setminus\{0\}$, define $F_a(x)=x^2-a^2$. At zero,
\[
F_a'(0)=0,
\]
and the correction equation $0\cdot u=a^2$ has no solution. Nevertheless, $F_a(a)=0$.
\end{proposition}
\begin{proof}
Differentiation gives $F_a'(x)=2x$, so the derivative at zero is zero. Since $a\ne0$, its square is nonzero, while $0\cdot u=0$ for every $u$. Finally, $F_a(a)=a^2-a^2=0$.
\end{proof}
\formal{Nonlinear.nonlinear_root_despite_first_order_obstruction}

The root can lie arbitrarily close to zero by choosing a small nonzero $a$. This produces a family of nearby counterexamples. It does not assert that one fixed equation with a nonzero residual at zero has roots converging to zero.

\begin{proposition}[A solvable linearization without a real root]\label{prop:no-real-root}
Let $G(x)=x^2+x+1$. Then
\[
G'(0)=1,\qquad \exists u\in\R\ (u=-1),
\]
so the linearized correction at zero is solvable. However, $G$ has no real root.
\end{proposition}
\begin{proof}
The derivative is $2x+1$, which equals one at zero, and $G(0)=1$. Thus $u=-1$ solves the linearized equation. But
\[
G(x)=\left(x+\frac12\right)^2+\frac34>0
\]
for every real $x$.
\end{proof}
\formal{Nonlinear.linearized_solution_without_real_root}

\section{Three conclusions that must stay distinct}
The results support three different kinds of certificate.

\textbf{An exact linear repair} provides witnesses for $Du+Ew=-r$ in the original target. Its validity is governed by quotient image membership.

\textbf{A finite arithmetic lift} provides a correction at one modulus and proves the exact divisibility gain. Its validity also uses arithmetic conditions such as B\'ezout coprimality and a scale divisibility relation.

\textbf{A nonlinear existence result} proves that a solution of the original nonlinear equation exists. It requires an additional theorem controlling the nonlinear terms, compatible successive approximations, or a different existence mechanism.

Paper III develops correction, nonlinear analysis, and filtered lifting under separate hypotheses~\cite{paperiii}. The present formal component does not import those complete releases. It proves a restricted linear theory, a finite arithmetic step, and counterexamples that prevent unjustified upgrades between these conclusions.

\section{The relation to Realization Calculus}\label{sec:realization}
The conceptual setting is the established separation between a mathematical object, its chosen realization, and the operations a realization preserves. Paper I's semantic realization criterion requires an explicit construction, preservation of the operations used in the proof, verification of additional domain hypotheses, and a conclusion stated in native objects~\cite{paperi}.

Constructive repair fits that criterion in a direct way. The native equation supplies $V$ and its residuals. The allowed transformations supply $D$ and the permissible class of added maps. The quotient organizes which residual classes those maps can remove. A transport that is used to decide the problem must preserve the necessary obstruction data.

Komorebi Space is the proposed ambient setting for organizing realizations, compatibility data, residuals, and obstructions. This monograph supplies a bounded repair construction that can be placed in that setting. It does not prove that every possible realization in Komorebi Space has a canonical finite obstruction space or an optimal repair.

A successful application must still identify the formal data with the actual objects of its domain. For example, an arithmetic representation with prescribed character and phase information cannot be obtained merely by naming an abstract linear space an arithmetic carrier. The Paper XI research companion distinguishes finite comparison from arithmetic existence and develops specific enlarged carrier constructions under stated hypotheses~\cite{paperxi}. Those constructions motivate the repair viewpoint, but are not reproduced or newly formalized here.

\section{Established reach and remaining questions}
The finite theory has a clear end point. It proves exactly when an added linear map repairs a family, the least possible added dimension over a field, optimal attainment inside a feasible legal subspace, complete finite menu feasibility, and a minimum cost search over finite domains. Its transport theorem excludes certificates that work only by collapsing the relevant quotient information.

The rational matrix bridge in Chapters~\ref{ch:rational} and~\ref{ch:matrix-attainment} now connects an exact Python construction to a small formal verifier. Further questions remain separate. Equivariant, filtered, topological, and arithmetic repairs need admissibility classes more structured than a subspace. Nonlinear minimum repair needs conditions that do not confuse a derivative obstruction with global nonexistence. Infinite lifting needs compatible choices and a convergence theorem.

These are extensions of a precise contract rather than reasons to enlarge its present claims. The elementary linear algebra and arithmetic mechanisms are classical. The contribution of this component is their explicit repair interface, formal proofs, certified finite search, and preservation boundaries. No claim of a universally complete discovery procedure or historical priority follows from the formal theorem count.

\begin{scopebox}[title=A usable stopping point]
A certificate must state which equation it solves, which changes are allowed, and which information is preserved. Within the declared finite linear setting, the theorems provide both successful repairs and complete rejection of the supplied menu. Outside that setting, the missing domain hypotheses remain mathematical obligations.
\end{scopebox}

\chapter{Exact rational certificates}\label{ch:rational}
\section{Let the search propose and the kernel check}
The finite optimizer in Chapter~\ref{ch:search} enumerates finite correction domains. Rational coefficients do not form such a domain. The appropriate interface is therefore not an enumeration of all rational assignments. A solver proposes finite rational matrices, and a verifier checks their identities in the original target.

Fix matrices
\[
 D\in\Q^{m\times n},\qquad R\in\Q^{m\times s},\qquad E\in\Q^{m\times q}.
\]
The residuals are the columns of $R$. For a selected menu $S$, let $\widetilde E_S$ retain the original candidate indexing and replace unselected columns by zero. The supplied coefficient matrices have shapes $U\in\Q^{n\times s}$ and $W\in\Q^{q\times s}$. Their correction equation is
\begin{equation}\label{eq:rational-original}
 DU+\widetilde E_SW=-R.
\end{equation}
Keeping the original indexing avoids any ambiguity about a permutation of candidate columns. Coefficients on a zeroed column have no effect. The admissibility predicate and cost function belong to the original problem, not to the certificate producer.

\begin{proposition}[Successful certificate]\label{prop:rational-success}
If an admissible $S$ and rational matrices $U,W$ satisfy \eqref{eq:rational-original}, they repair every original residual.
\end{proposition}
\begin{proof}
Take the $i$th column of the displayed identity. It is the original equation $Du_i+\sum_{j\in S}w_{ji}e_j=-r_i$.
\end{proof}
\formal{RationalCertificates.checkRepair_sound}

\begin{scopebox}[title=A rejected certificate is not a rejected problem]
Failure to parse data, a false identity, a missing witness, or a resource limit is an unresolved verification outcome. It does not establish mathematical infeasibility. A rejection of the problem requires a valid obstruction certificate with an explicit scope.
\end{scopebox}

\section{Separating obstructions}
\begin{definition}[Separating row]
A separating certificate for $S$ is a row $Y\in\Q^{1\times m}$ satisfying
\begin{equation}\label{eq:rational-separator}
 YD=0,\qquad Y\widetilde E_S=0,\qquad YR\ne0.
\end{equation}
The last condition asks that at least one original residual is detected. No residual index needs to be selected when the family is empty.
\end{definition}

\begin{theorem}[Sound separation]\label{thm:rational-separation}
A separating certificate rules out a repair by the selected menu.
\end{theorem}
\begin{proof}
Multiplying a supposed correction equation on the left by $Y$ gives $0=-YR$, contradicting \eqref{eq:rational-separator}.
\end{proof}
\formal{RationalCertificates.annihilator_kills_residual}
\formal{RationalCertificates.checkSeparator_sound}

\begin{proposition}[Linear alternative]\label{prop:rational-alternative}
For a fixed rational menu, either the correction matrices exist or a separating row exists, but not both.
\end{proposition}
\begin{proof}
Exclusivity follows from the preceding theorem. If no correction matrices exist, some residual $r_i$ lies outside the column space of $[D\ \widetilde E_S]$. Choose a basis for that space, adjoin $r_i$, and extend to a basis of $\Q^m$. The functional taking value one on $r_i$ and zero on all other chosen basis vectors is a rational separating row.
\end{proof}
\sourcescope{This completeness argument is an ordinary proof. The new Lean module proves certificate soundness, not a general algorithm or formal theorem producing a separator for every infeasible rational menu.}

\section{Dimension certificates without a reported rank}
Suppose an augmentation with $t$ columns succeeds. Its optimality can be certified by two additional matrices
\[
 H\in\Q^{t\times m},\qquad Q\in\Q^{s\times t}
\]
with
\begin{equation}\label{eq:rational-dimension}
 HD=0,\qquad HRQ=I_t.
\end{equation}
These identities make a lower bound on the required dimension explicit.

\begin{theorem}[Dimension lower certificate]\label{thm:rational-lower}
Assume \eqref{eq:rational-dimension}. If an arbitrary competing augmentation $E'\in\Q^{m\times w}$ has coefficients satisfying $DU'+E'W'=-R$, then $t\leq w$.
\end{theorem}
\begin{proof}
Multiplication by $H$ and $Q$ gives $HE'W'Q=-I_t$, hence
\[
 (HE')(-W'Q)=I_t.
\]
Thus $HE':\Q^w\to\Q^t$ has a right inverse. It is surjective, and consequently $t\leq w$.
\end{proof}
\formal{RationalCertificates.dimension_certificate_right_inverse}
\formal{RationalCertificates.dimension_certificate_lower_bound}

A successful augmentation with $t$ columns attains this bound. The verifier checks the original repair identity and \eqref{eq:rational-dimension}. It does not trust a numerical rank supplied by the search program.
\formal{RationalCertificates.checkDimension_sound}

\begin{proposition}[Availability of lower certificates]
Let $O$ be the span of the residual classes in $\Q^m/\im D$, and let $t=\dim O$. Then matrices $H,Q$ satisfying \eqref{eq:rational-dimension} exist.
\end{proposition}
\begin{proof}
Choose coordinates $O\cong\Q^t$, extend that map to the whole quotient, and compose with the quotient map to obtain $H$. Then $HD=0$. The columns of $HR$ span $\Q^t$, so $HR$ has a linear right inverse $Q$.
\end{proof}
\sourcescope{The availability argument is an ordinary finite-dimensional proof. The concrete producer in the next chapter finds and exports such data. Its search implementation is not itself proved correct in Lean.}

\section{Exact rank and menu cost}
For $A\in\Q^{a\times b}$, a certificate of rank $d$ consists of appropriately shaped rational matrices with
\begin{equation}\label{eq:rational-rank}
 A=FG,\qquad PAQ=I_d,
\end{equation}
where the inner dimension of $FG$ is $d$.

\begin{proposition}[Rank certificate]
The identities \eqref{eq:rational-rank} imply $\rank A=d$.
\end{proposition}
\begin{proof}
The factorization gives $\rank A\leq d$. The identity factorization through $A$ gives $d\leq\rank A$.
\end{proof}
\formal{RationalCertificates.checkRank_sound}

For comparison with the original quotient statement, ordinary rank-nullity gives
\begin{equation}\label{eq:rank-difference-new}
 \dim O=\rank[D\ R]-\rank D.
\end{equation}
Indeed, the quotient map restricted to $\im D+\im R$ has kernel $\im D$ and image $O$. The current matrix verifier does not depend on a formal identification of this formula with the earlier abstract Lean quotient API. It obtains the lower bound directly from a right inverse.

Minimum menu cost is a separate optimization problem. To certify a successful menu $S$, the verifier requires a separating row for every admissible menu $T$ with strictly smaller declared cost. It quantifies over all subsets of the original finite candidate universe. It does not trust an incomplete producer-supplied list.

\begin{theorem}[Certified finite-menu optimum]
An admissible successful menu, together with separating rows for every cheaper admissible menu, has minimum cost among all admissible feasible menus.
\end{theorem}
\begin{proof}
A cheaper feasible menu would contradict its supplied separating row by Theorem~\ref{thm:rational-separation}.
\end{proof}
\formal{RationalCertificates.checkOptimal_sound}

A certificate rejecting every admissible subset proves that no menu in the declared universe succeeds.
\formal{RationalCertificates.checkNoRepair_sound}

The rational example $D(u)=(u,0,0)$, $r=(0,1,0)$, $e_0=(0,1,1)$, $e_1=(0,0,1)$ is checked in the new matrix layer. Both entries repair the residual with coefficients $(-1,1)$. Rows $(0,1,0)$, $(0,1,-1)$, and $(0,1,0)$ separate the empty menu, the first singleton, and the second singleton. Thus the minimum cardinality is two. Restricting admissibility to at most one entry gives a checked complete rejection.
\formal{RationalMenus.menu_certificate}
\formal{RationalMenus.two_entries_minimal}
\formal{RationalMenus.restricted_certificate}
\formal{RationalMenus.no_small_menu}

The subset search may be exponential in menu size. The free augmentation exporter developed next does not implement a scalable menu optimizer, and the existence of a free minimum says nothing about which directions occur in a prescribed menu.

\chapter{Constructive matrix attainment}\label{ch:matrix-attainment}
\section{A complete detector}
Fix $D\in\Q^{m\times n}$ and $R\in\Q^{m\times s}$. A split correction certificate consists of
\[
 P\in\Q^{n\times m},\quad K\in\Q^{m\times d},\quad L\in\Q^{d\times m}
\]
with
\begin{equation}\label{eq:split-certified}
 LD=0,\qquad DP+KL=I_m.
\end{equation}
The first identity makes $L$ a detector that ignores original corrections. The second proves that no further information has been discarded: every target vector has the decomposition $v=DPv+KLv$. There is no assumption that $LK=I_d$, and redundant detector coordinates are permitted.

\begin{theorem}[Exact detection]\label{thm:split-exact}
Under \eqref{eq:split-certified},
\[
 (\exists X,\ DX=-R)\quad\Longleftrightarrow\quad LR=0.
\]
When $LR=0$, the explicit witness is $X=-PR$.
\end{theorem}
\begin{proof}
A correction gives $0=LDX=-LR$. Conversely, $LR=0$ gives $R=(DP+KL)R=DPR$, so $D(-PR)=-R$.
\end{proof}
\formal{RationalCertificates.split_correctable_iff}

At the quotient level, $L$ induces an injection $\Q^m/\im D\to\Q^d$. If $Lv=0$, the split identity gives $v=DPv$, proving that its class is zero. The image is $\im L$. Hence the required residual obstruction span is isomorphic to the column space of $C=LR$. This quotient interpretation is an ordinary consequence of the displayed identities.

\section{The attaining construction}
Suppose the detected residual matrix $C=LR$ has certificate data
\[
 F\in\Q^{d\times t},\quad G\in\Q^{t\times s},\quad
 J\in\Q^{t\times d},\quad Q\in\Q^{s\times t}
\]
satisfying
\begin{equation}\label{eq:factor-certified}
 LR=FG,\qquad JLRQ=I_t.
\end{equation}
The new augmentation and its correction coefficients are not unspecified witnesses. They are
\begin{equation}\label{eq:attaining-formula}
 \boxed{E=KF,\qquad U=-PR,\qquad W=-G.}
\end{equation}

\begin{theorem}[Constructive minimum augmentation]\label{thm:matrix-attainment}
Assume \eqref{eq:split-certified} and \eqref{eq:factor-certified}. The matrices in \eqref{eq:attaining-formula} satisfy $DU+EW=-R$. Every competing rational augmentation repairing the same residual family has at least $t$ added dimensions.
\end{theorem}
\begin{proof}
The constructed matrices satisfy
\[
 DU+EW=-DPR-KFG=-DPR-KLR=-(DP+KL)R=-R.
\]
For the lower bound set $H=JL$. Then $HD=0$ and $HRQ=I_t$. Apply Theorem~\ref{thm:rational-lower} to every competing repair.
\end{proof}
\formal{RationalCertificates.split_constructs_repair}
\formal{RationalCertificates.checkOptimalConstruction_sound}

This proves attainment and minimality after checking four finite identities. The generalized inverse and the choice of a residual basis belong to the producer. They are not assumed trustworthy. A valid set of identities is sufficient, regardless of how the producer found it.

\section{How the Python producer obtains the matrices}
The executable producer uses exact rational row reduction. It selects independent columns of $D$ and extends them by standard basis vectors to a basis $B$ of $\Q^m$. The inverse $B^{-1}$ supplies the image coordinates and the complementary detector coordinates. The selected image coordinates are placed in the original domain to obtain $P$, while the complementary basis vectors form $K$ and the corresponding rows of $B^{-1}$ form $L$.

For $C=LR$, the producer selects independent residual columns to form $F$. A selection matrix $Q$ records those columns. A left inverse $J$ of $F$ gives $G=JC$. The producer then forms $E,U,W$ using \eqref{eq:attaining-formula}, and checks every identity again with exact rational arithmetic before export.

All of these data can be obtained over a field. The search implementation may nevertheless contain a bug, encounter a resource limit, or reject malformed input. Such events do not become mathematical no-go results. The Lean certificates verify successful outputs, not the total correctness or complexity of this search program.

Zero-dimensional targets, original domains, obstruction spaces, and residual families are retained. For an already correctable residual family the construction has $t=0$, so it adds no directions. The input limit of the supplied prototype is an implementation resource bound, not a restriction on Theorem~\ref{thm:matrix-attainment}.

\section{A full intermediate trace}\label{sec:matrix-worked}
Take
\[
 D=\begin{pmatrix}1\\0\\0\end{pmatrix},\qquad
 R=\begin{pmatrix}2&4\\1&2\\1&2\end{pmatrix}.
\]
The producer selects the first column of $D$ and the complementary second and third coordinate axes. It obtains
\[
 P=\begin{pmatrix}1&0&0\end{pmatrix},\quad
 K=\begin{pmatrix}0&0\\1&0\\0&1\end{pmatrix},\quad
 L=\begin{pmatrix}0&1&0\\0&0&1\end{pmatrix}.
\]
Thus $LD=0$, $DP=\operatorname{diag}(1,0,0)$, and $KL=\operatorname{diag}(0,1,1)$. The residual coordinates are
\[
 C=LR=\begin{pmatrix}1&2\\1&2\end{pmatrix}.
\]
A rank-one certificate is
\[
 F=\begin{pmatrix}1\\1\end{pmatrix},\quad
 G=\begin{pmatrix}1&2\end{pmatrix},\quad
 J=\begin{pmatrix}1&0\end{pmatrix},\quad
 Q=\begin{pmatrix}1\\0\end{pmatrix}.
\]
The result is
\[
 E=\begin{pmatrix}0\\1\\1\end{pmatrix},\qquad
 U=\begin{pmatrix}-2&-4\end{pmatrix},\qquad
 W=\begin{pmatrix}-1&-2\end{pmatrix}.
\]
The original equations become
\[
 DU+EW=
 \begin{pmatrix}-2&-4\\0&0\\0&0\end{pmatrix}+
 \begin{pmatrix}0&0\\-1&-2\\-1&-2\end{pmatrix}=-R.
\]
The lower witness is $H=JL=(0,1,0)$, with $HD=0$ and $HRQ=1$. The ambient cokernel has dimension two, but the residual classes span only one direction. The construction therefore adds one direction and proves that zero added directions are insufficient.
\formal{Case000.certificate_checked}
\formal{Case000.original_equations_checked}
\formal{Case000.optimal}

The fractional case uses $D=(2/3,4/3)^{\mathsf T}$ and $R=(1/2,5/6)^{\mathsf T}$. It exercises nonintegral splitting and correction coefficients while retaining exact arithmetic. Empty and rank-deficient cases are included in the same exported corpus. Every case has its own original input file, certificate file, and formal optimality statement.

\section{Serialization and input identity}
The wire format represents a rational entry by a signed decimal numerator and a positive decimal denominator. The Python decoder rejects floats, nonfinite values, duplicate object keys, zero or negative denominators, unreduced pairs, malformed dimensions, and extraneous certificate fields. Resource limits produce an invalid or unresolved outcome, not an infeasibility theorem.

The generated Lean data use typed pairs in $\Z\times\N$. Their decoder returns the intended rational when the denominator is nonzero and explicitly rejects a zero denominator at the cell level. Each generated input and certificate additionally proves denominator validity for all its cells.
\formal{RationalCertificates.decodeCell_zero}
\formal{RationalCertificates.decodeCell_nonzero}
\formal{RationalCertificates.decodeCell_sound}
\formal{RationalCertificates.checkWireMatrix_sound}
\formal{RationalCertificates.decodeWireMatrix_cell}

The original matrices $D,R$ are emitted separately from certificate data. The certificate module imports them and cannot replace them through a certificate field. A canonical problem hash binds the JSON records. A diagnostic Lean echo of the decoded input matrices is compared with the original JSON values for every exported case. That echo is an additional execution check, never a theorem premise.

\begin{scopebox}[title=The precise trust boundary]
Lean checks the matrix identities, the original equations, and the minimum-dimension conclusion for the literal decoded problem in the generated input module. The numeric cell-decoding lemmas are formalized. The whole JSON parser, file system, hashing routine, and source emitter are not proved correct in Lean. Their input binding is checked by hashes, strict decoding, and a separate Lean echo. A user must still identify the intended original problem correctly.

No theorem is accepted through \code{native_decide}. Concrete certificate propositions are proved with \code{decide +kernel}, and every exported theorem is included in the transitive axiom audit.
\end{scopebox}

\section{Legality and coefficient rings}
Free minimum dimension is not minimum menu cost. The constructed direction $KF$ need not occur among supplied menu entries. When legal directions form $\im A$, an additional identity $KF=AZ$ certifies that the augmentation lies in that legal subspace. This identity also rewrites every added correction as $A(ZW)$.
\formal{RationalCertificates.constructed_repair_in_legal_image}

The availability of splitting also depends on the coefficient ring. For $D=(2)$ over $\Z$, the condition $LD=0$ forces $L=0$. The split identity would then require $2P=1$, which is impossible over $\Z$. Over $\Q$, $P=1/2$ is permitted. The current producer and checker are rational, not a certificate of integral realizability.

Finally, exact rational correction is not nonlinear convergence or arithmetic representation-theoretic existence. The existing boundaries of Paper III and Conjecture 66 are unchanged. The contribution of this extension is a usable and checked witness interface inside the established calculus, not a replacement for its foundations.

\chapter{Costed rational menus and rejection covers}
\label{ch:menu-covers}

The free augmentation problem asks how many new dimensions are necessary. A prescribed menu asks a different question: which available directions should be purchased, enabled, or admitted? Two menu entries can be needed for a one-dimensional obstruction, and the cheapest solution need not use the fewest entries. This chapter makes that second problem operational over the rationals.

The implementation uses exact elimination to propose an answer. A separate certificate checker verifies the answer against the original matrices and the complete admissibility rule. The search procedure is outside the proof trust boundary. The returned mathematical witnesses are not.

\section{The original menu contract}

Fix rational matrices
\[
 D\in\Q^{m\times n},\qquad E\in\Q^{m\times q},\qquad R\in\Q^{m\times s}.
\]
The columns $e_j$ of $E$ are the available new directions. The finite family
\[
 \mathcal A\subseteq\mathcal P(\{0,\ldots,q-1\})
\]
contains exactly the admissible selections. It need not be downward closed. Nonnegative integer weights $c_j$ define the cost
\[
 c(S)=\sum_{j\in S}c_j.
\]
These matrices, selections, and costs are fixed original input. A certificate cannot replace them with another problem.

Let $\widetilde E_S$ retain each column of $E$ selected by $S$ and replace every other column by zero. This convention preserves the original candidate indexing. A coefficient table $W\in\Q^{q\times s}$ may have arbitrary values in unselected rows, but those rows contribute nothing. A proposed repair consists of matrices $U,W$ with
\begin{equation}
 DU+\widetilde E_SW=-R.
 \label{eq:menu-cover-primal}
\end{equation}
Every selected direction is charged, even if its coefficient row happens to be zero. Thus cost measures the selected menu, not merely the support of a particular coefficient table.

\section{Separating blocks}

An infeasible menu admits a rational separating row, by the exact linear alternative in \cref{ch:rational}. The row annihilates all available corrections while detecting at least one residual. Instead of providing a separate row for every rejected subset, we can reuse one row on all subsets of a larger rejected set.

\begin{definition}[Separating block]
A separating block is a pair $(B,Y)$, where $B\subseteq\{0,\ldots,q-1\}$ and $Y\in\Q^{1\times m}$ satisfy
\begin{equation}
 YD=0,\qquad Y\widetilde E_B=0,\qquad YR\ne0.
 \label{eq:menu-cover-dual}
\end{equation}
The set $B$ need not itself belong to $\mathcal A$.
\end{definition}

The last condition is a matrix inequality: at least one column is detected. For an empty residual family it cannot hold, which correctly prevents a false separating witness for a vacuously solved family.

\begin{lemma}[Column-mask identity]
For every row $Y$ and column $j$,
\[
 (Y\widetilde E_S)_{0j}=
 \begin{cases}
 (YE)_{0j},&j\in S,\\
 0,&j\notin S.
 \end{cases}
\]
\end{lemma}
\begin{proof}
The mask either retains the entire $j$th column or replaces it by zero. Expanding the finite matrix product gives the two cases.
\end{proof}
\formal{KomorebiRepair.RationalCertificates.mul_masked_entry}

\begin{theorem}[Downward validity of a separator]
If $(B,Y)$ is a separating block and $T\subseteq B$, then $Y$ also separates $T$. In particular, $T$ cannot satisfy \eqref{eq:menu-cover-primal}.
\end{theorem}
\begin{proof}
For a selected column $j\in T$, one also has $j\in B$, so $(YE)_{0j}=0$ by the block identity. Every unselected column is zero after masking. Thus $Y\widetilde E_T=0$. The other two equations in \eqref{eq:menu-cover-dual} are unchanged.

If a repair for $T$ existed, multiplication by $Y$ would give
\[
 0=YDU+Y\widetilde E_TW=-YR,
\]
contradicting the detection condition.
\end{proof}
\formal{KomorebiRepair.RationalCertificates.separates_subset}

The restriction to an admissible family does not invalidate this argument. A block can be a forbidden superset while still supplying valid mathematical evidence against its permitted subsets. What matters is that the block's own linear impossibility has been checked.

\section{Complete coverage of cheaper alternatives}

Let $\mathcal B$ be a finite list of separating blocks. Say that $\mathcal B$ covers $T$ when some $(B,Y)\in\mathcal B$ satisfies $T\subseteq B$. Validity of the block list and coverage of the relevant selections are separate obligations.

\begin{theorem}[Certified minimum menu cost]
Suppose the following data have been checked:
\begin{enumerate}
 \item $S\in\mathcal A$ and the original repair equation holds for $S,U,W$.
 \item Every member of $\mathcal B$ is a valid separating block.
 \item Every $T\in\mathcal A$ with $c(T)<c(S)$ is covered by $\mathcal B$.
\end{enumerate}
Then $S$ has minimum cost among all admissible feasible selections.
\end{theorem}
\begin{proof}
The first condition supplies an admissible feasible selection. If a cheaper admissible feasible $T$ existed, the third condition would place it inside a separating block. Downward validity would make $T$ infeasible, a contradiction.
\end{proof}
\formal{KomorebiRepair.RationalCertificates.checkCoveredOptimal_sound}

The checker quantifies over the whole finite candidate-subset universe when it verifies coverage. It does not accept a producer's claim that its own list of alternatives is complete. Equal-cost alternatives need not be rejected, and the certificate does not claim that the optimum is unique.

\begin{theorem}[Certified rejection of the admissible family]
If every block is valid and every member of $\mathcal A$ is covered, then no admissible menu repairs the residual family.
\end{theorem}
\begin{proof}
Every admitted selection lies inside a valid separating block and is therefore infeasible.
\end{proof}
\formal{KomorebiRepair.RationalCertificates.checkCoveredNoRepair_sound}

When $\mathcal A$ is empty, this conclusion is vacuous. The implementation reports that fact explicitly. It does not manufacture a nonzero obstruction or claim that the original equations have no solution without the admissibility restriction.

\begin{scopebox}[title=The scope of rejection]
An accepted rejection certificate rules out the declared admissible menus. It does not rule out a new direction outside the menu, a different coefficient ring, a changed admissibility rule, or a different arithmetic host. A failed check or exhausted time budget establishes none of these impossibility claims.
\end{scopebox}

\section{Constructing the evidence}

For each candidate subset, the producer performs exact rational row reduction on the selected correction columns. It keeps the row transformation matrix as part of the computation. If the transformed residual has a nonzero entry in a zero coefficient row, that row transformation supplies a separating functional. Otherwise the pivot equations give a coefficient table for every residual. Each proposed identity is rechecked using exact arithmetic.

The producer evaluates all $2^q$ selections in this bounded implementation. It chooses a feasible admitted selection of minimum declared cost, with a deterministic tie break used only for reproducibility. Alternatively, it records that none of the admitted selections passed the exact linear alternative.

For rejection evidence, the producer finds inclusion-maximal infeasible subsets of the full candidate universe. Every infeasible selection is contained in such a maximal subset because the universe is finite. A separating row for each selected maximal subset therefore provides a complete supply of potential rejection blocks. The producer then greedily retains blocks that cover the required alternatives.

The greedy compression is not claimed to minimize certificate size. Nor does a smaller rejection cover imply a faster search algorithm: the present producer still examines every subset. The verified conclusion concerns menu cost and the original equations, not optimal compression or polynomial-time complexity.

\section{Two entries for one missing dimension}

Take
\[
 D=\begin{pmatrix}1\\0\\0\end{pmatrix},\qquad
 E=\begin{pmatrix}0&0\\1&0\\1&1\end{pmatrix},\qquad
 R=\begin{pmatrix}0\\1\\0\end{pmatrix},
\]
with both candidate weights equal to one and every subset admitted. The chosen menu is $S=\{0,1\}$, with
\[
 U=(0),\qquad W=\begin{pmatrix}-1\\1\end{pmatrix}.
\]
Its exact repair is
\[
 DU+EW=\begin{pmatrix}0\\-1\\0\end{pmatrix}=-R.
\]
Two separating blocks suffice:
\[
 B_0=\{0\},\quad Y_0=(0,1,-1),\qquad
 B_1=\{1\},\quad Y_1=(0,1,0).
\]
Both satisfy $Y_iD=0$, annihilate their masked candidate matrices, and give $Y_iR=1$.

\begin{center}
\begin{tabular}{@{}ccc@{}}\toprule
Cheaper menu & Cost & Covering block\\\midrule
$\varnothing$ & 0 & $B_0$ or $B_1$\\
$\{0\}$ & 1 & $B_0$\\
$\{1\}$ & 1 & $B_1$\\\bottomrule
\end{tabular}
\end{center}
Thus cost two is minimal. Restricting admissibility to the three cheaper menus turns the same two blocks into a complete no-admissible-repair certificate. The unrestricted one-dimensional repair direction $R$ still exists, but it is not one of the two listed candidate directions.

\sourcescope{Executed examples: Menu000 and Menu001. Each has separately emitted original input, exact rational witnesses, a Lean certificate proof, an axiom audit, and an input echo.}

\section{Cost, dimension, and evidence size}

A third candidate can make the cost distinction sharper. Add the direct direction $e_0=(0,1,0)$ at cost five, and price $e_1=(0,1,1)$ and $e_2=(0,0,1)$ at one each. One menu entry can now repair the residual, but the cheapest repair selects two entries, $\{1,2\}$, at total cost two. Minimum free dimension, minimum menu cardinality, and minimum menu cost are three different quantities.

A separate example shows why rejection covers help. Let
\[
 D=\begin{pmatrix}1\\0\end{pmatrix},\qquad
 E=\begin{pmatrix}2&3&4&5\\0&0&0&0\end{pmatrix},\qquad
 R=\begin{pmatrix}0\\1\end{pmatrix}.
\]
The row $Y=(0,1)$ separates the full four-entry menu. One block therefore excludes all sixteen subsets. Adding $(0,1)$ as a new candidate would change the problem and allow a repair. The certified rejection of the original menu would remain correct.

\sourcescope{Executed examples: Menu002 distinguishes cost from cardinality. Menu004 excludes all sixteen subsets with one separating block. Menu005 uses one block to cover eight cheaper selections before accepting a positive-cost direction.}

\section{Fractional witnesses and independent original data}

The solver is not restricted to coefficients in a finite field. For
\[
 D=\begin{pmatrix}1\\0\end{pmatrix},\quad
 E=\begin{pmatrix}0\\2/3\end{pmatrix},\quad
 R=\begin{pmatrix}1/2\\5/6\end{pmatrix},
\]
the selected direction at cost seven has witnesses
\[
 U=-1/2,\qquad W=-5/4.
\]
The checker verifies the exact rational equation, with no tolerance or floating-point rank test.

The generated input module contains $D,E,R$, the candidate weights, and the literal admissibility table. The certificate module imports those inputs and supplies only coefficient matrices and separating blocks. Typed rational decoding and nonzero denominator proofs are retained from the preceding bridge.

An independent diagnostic echo evaluates the matrices in Lean and evaluates the admissibility and cost of every candidate subset. Python compares those evaluated results with the original JSON input. This comparison catches mismatched original constraints as well as mismatched matrices. It is an executed statement-binding check, not a formal proof of the entire external JSON parser or source emitter.

\section{Bounded executable use}

The command \code{scripts/certify_menu.py} creates a new local directory, constructs or reads a certificate, emits the fixed input and witness modules, compiles their proofs, audits their axioms, replays the modules, and checks the original-data echo. Existing output directories are refused.

\begin{Verbatim}[fontsize=\small,breaklines=true]
export KOMOREBI_LEAN_KIT=/path/to/Komorebi_Portable_Lean_4.34.0-rc2
python scripts/certify_menu.py examples/menu_corpus/Menu003_input.json --output /tmp/repair-menu
\end{Verbatim}

The successful mathematical outcomes are \nolinkurl{KERNEL_CHECKED_OPTIMAL_MENU} and \nolinkurl{KERNEL_CHECKED_NO_ADMISSIBLE_MENU}. Invalid data, unsupported sizes, compiler failure, and timeouts produce \code{UNRESOLVED}. A Python-valid certificate alone is not a kernel-checked result.

This version limits the menu to eight candidates. Other input bounds include thirty-two rows or columns, canonical rational cells with bounded decimal length, and a four-megabyte JSON file limit. These are implementation safeguards, not hypotheses of the finite mathematical theorems. Exhaustive search is still exponential in menu size, and the time limit can be reached within the admitted size bounds.

The module is a constructive extension of the correction framework already developed in Realization Calculus. It gives Komorebi Space a more explicit candidate interface for admissible finite linear repairs, but supplies no universal semantic-realization theorem. Integration into the authoritative Paper III tree and the established Komorebi Python library remains separate from this isolated, reproducible implementation.

\chapter{Proof carrying branch search}
\label{ch:branch-search}

The preceding menu verifier checks every required selection against a family of separating blocks. That is a sound endpoint, but its coverage test still ranges over the full subset universe. A large menu may therefore be expensive to verify even when one short argument excludes almost every alternative.

This chapter replaces the powerset coverage test with a finite search tree. Each node describes an interval of selections. A terminal node either proves a cost lower bound, supplies one separator for the entire interval, or checks that the original admissibility rule excludes that interval. A branching node must retain both possible decisions about one unresolved candidate.

The producer may search in any order and use any heuristic. The verifier checks the submitted tree. It does not trust the producer's account of which alternatives were examined.

\begin{scopebox}[title=The new verified boundary]
The mathematical result is soundness of a finite rational certificate tree. It certifies minimum cost within the declared menu or infeasibility of the declared admissible family. It is not a proof of correctness or termination of the entire Python producer, a polynomial complexity result, or a native arithmetic realization theorem.

All 37 earlier proof modules are retained unchanged. The new module reuses their exact rational repair and separator theorems.
\end{scopebox}

\section{Intervals of candidate selections}

Let $J=\{0,\ldots,q-1\}$ be the original candidate index set. The matrices $D,E,R$ and the original correction equation remain fixed:
\[
 DU+\widetilde E_S W=-R.
\]
Here $\widetilde E_S$ zeros unselected columns of $E$ while retaining its original indexing. Costs are nonnegative integers attached to those same columns:
\[
 c(S)=\sum_{j\in S}w_j,\qquad w_j\in\N.
\]
The admissibility rule either permits every subset or supplies a fixed finite table of permitted subsets. An empty table permits none. It is different from the all-subsets rule.

\begin{definition}[Selection interval]
For $A\subseteq B\subseteq J$, define
\[
 \mathcal I(A,B)=\{T\subseteq J:A\subseteq T\subseteq B\}.
\]
The set $A$ contains candidates already required, and $B$ contains all candidates still permitted by the branch decisions. The unresolved indices form $B\setminus A$.
\end{definition}

\begin{lemma}[Exact branch partition]
If $j\in B\setminus A$, then
\[
 \mathcal I(A,B)=
 \mathcal I(A,B\setminus\{j\})\;\sqcup\;
 \mathcal I(A\cup\{j\},B).
\]
\end{lemma}
\begin{proof}
Every $T\in\mathcal I(A,B)$ either excludes $j$ or includes $j$. In the first case it belongs to the left interval, and in the second it belongs to the right interval. Conversely, each child interval is contained in the parent. Their membership conditions on $j$ are incompatible, so the union is disjoint.
\end{proof}

This identity is the completeness mechanism of the verifier. A certificate cannot simply omit an inconvenient branch. Its split constructor contains both children, and the checker verifies both. The proof of \decl{checkBranch_sound} applies the two cases directly.

\section{Three terminal certificates}

\subsection{Cost lower bound}

Nonnegative additive costs are monotone under inclusion:
\[
 A\subseteq T\quad\Longrightarrow\quad c(A)\leq c(T).
\]
Therefore, when an incumbent repair has cost $c_\ast$, the finite check
\[
 c_\ast\leq c(A)
\]
excludes every strictly cheaper repair in $\mathcal I(A,B)$. It does not say the interval is infeasible. It only says that it cannot improve the incumbent.
\formal{KomorebiRepair.RationalCertificates.additiveCost_mono}

\begin{proof}
The sum for $c(T)$ contains all terms in the sum for $c(A)$ and only additional nonnegative terms. Apply transitivity to $c_\ast\leq c(A)\leq c(T)$.
\end{proof}

A cost leaf is forbidden in an impossibility certificate. Even a very expensive feasible repair refutes a claim that no repair exists. The implementation expresses this distinction using an optional threshold: a supplied natural number requests a lower bound, while no threshold requests infeasibility.

\subsection{Separator for the upper selection}

A row $Y$ is checked against the largest selection $B$ in the interval:
\[
 YD=0,\qquad Y\widetilde E_B=0,\qquad YR\ne0.
\]
Every $T\subseteq B$ uses only columns annihilated by $Y$. Consequently, the same row proves that no $T\in\mathcal I(A,B)$ repairs the family. This reuses the downward separator theorem from the previous chapter.

The residual matrix in the last inequality is the original $R$. Projecting it to zero in some unrelated representation is not accepted as a repair.

\subsection{Exclusion by the original admissibility rule}

If admissibility is specified by a list $\mathcal A$, a terminal check establishes
\[
 \nexists T\in\mathcal A:\quad A\subseteq T\subseteq B.
\]
It examines the original list, not the powerset of $J$. For the all-subsets rule this terminal certificate is unavailable.
\formal{KomorebiRepair.RationalCertificates.intervalExcluded_sound}

It is not enough that $A$ itself be inadmissible. A larger selection in the interval might be permitted. For example, a rule that permits only the pair $\{0,1\}$ excludes the empty selection but not the interval from the empty selection to the full pair. The verifier retains this distinction and does not assume that admissibility is downward closed.

\section{Soundness of the whole tree}

\begin{definition}[Accepted branch tree]
A tree for $\mathcal I(A,B)$ consists of one of the three terminal certificates above, or a split at $j\in B\setminus A$ with accepted trees for both child intervals. In lower-bound mode a cost threshold is fixed throughout the tree. In infeasibility mode cost leaves are rejected.
\end{definition}

\begin{theorem}[Structural certificate soundness]
An accepted lower-bound tree with threshold $c_\ast$ proves
\[
 T\in\mathcal I(A,B),\quad T\text{ admissible and feasible}
 \quad\Longrightarrow\quad c_\ast\leq c(T).
\]
An accepted infeasibility tree proves that no admissible selection in $\mathcal I(A,B)$ is feasible.
\end{theorem}
\formal{KomorebiRepair.RationalCertificates.checkBranch_sound}

\begin{proof}
Induct on the supplied tree. For a cost leaf, monotonicity gives the lower bound. For a separator leaf, the exact linear alternative rules out feasibility of every subset of the upper selection. For an admissibility leaf, the original table rules out the premise.

At a split, the exact branch partition puts any proposed selection in one child interval. Apply that child's induction hypothesis. Both children are present in the data and checked by the verifier. This proves the assertion without enumerating every selection in the parent interval.
\end{proof}

\begin{corollary}[Minimum cost certificate]
Suppose an admissible selection $S$ has supplied coefficients $U,W$ satisfying the original repair equation. An accepted lower-bound tree of threshold $c(S)$ on $\mathcal I(\varnothing,J)$ proves that $S$ has minimum cost among all admissible feasible selections.
\end{corollary}
\formal{KomorebiRepair.RationalCertificates.checkTreeOptimal_sound}
\begin{proof}
The supplied coefficient equation proves feasibility of $S$. Every selection lies in $\mathcal I(\varnothing,J)$. Structural soundness therefore gives $c(S)\leq c(T)$ for every admissible feasible competitor $T$.
\end{proof}

\begin{corollary}[Scoped complete rejection]
An accepted infeasibility tree on $\mathcal I(\varnothing,J)$ proves that the declared admissible family contains no repair.
\end{corollary}
\formal{KomorebiRepair.RationalCertificates.checkTreeNoRepair_sound}

There is no claim that directions outside $J$ cannot repair the residual. Nor does a rejected certificate establish infeasibility. A missing child, a bad separator, or an exhausted resource budget leaves the outcome unresolved.

\section{How the producer constructs evidence}

The Python producer traverses intervals and retains exact row-reduction witnesses. It first checks whether the original admissibility table excludes the interval. It then uses the incumbent cost to test a bound leaf. If necessary, it solves the rational linear problem on the upper selection $B$. An infeasible upper selection supplies a separator that ends the whole branch.

If the lower selection $A$ is admitted and feasible, its coefficient matrix supplies a new incumbent and the interval ends at a cost leaf. Otherwise, the producer selects an unresolved index and explores both decisions. The current heuristic visits low-cost candidates first. Correctness of this heuristic is not an assumption of any theorem.

The incumbent may improve while the search proceeds. Earlier cost leaves remain sound: a later incumbent has no greater cost, so any earlier lower bound also excludes improvement over the final incumbent. The finished checker uses only the final threshold and rechecks every leaf against it. It need not trust the producer's history of incumbents.

Exact linear solves are cached by selection. This avoids repeating an identical solve when several intervals have the same upper or lower selection. Caching is an implementation choice outside the trusted mathematical layer.

\section{A twenty-four-candidate construction}

Let $b_1,b_2,b_3$ be the standard basis of $\Q^3$. Take $D(u)=ub_1$ and residual $r=b_2+b_3$. There are 24 candidate entries:
\[
 a_0=b_2,\qquad a_1=b_3,\qquad a_j=b_1\quad(2\leq j<24).
\]
Their costs are two, three, and ten for each remaining entry. All subsets are admissible.

The selected repair is $S=\{0,1\}$, with cost five, original coefficient zero, and coefficient $-1$ on each selected menu direction. Its full coefficient matrix satisfies the original equation exactly.

The proof tree has five nodes. At the root it splits on candidate zero. The excluded child lacks the only second-coordinate correction and has separator $(0,1,0)$. The included child splits on candidate one. Excluding candidate one has separator $(0,0,1)$. Including both reaches lower selection $\{0,1\}$ with cost five and terminates at a bound leaf.

\begin{center}
\begin{tabularx}{\linewidth}{@{}lX@{}}
\toprule Branch & Certified conclusion\\\midrule
Exclude candidate 0 & The second coordinate of the residual cannot be corrected.\\
Include 0, exclude 1 & The third coordinate of the residual cannot be corrected.\\
Include 0 and 1 & Every superset costs at least five.\\
\bottomrule
\end{tabularx}
\end{center}

There are $2^{24}=16{,}777{,}216$ possible selections, but the producer used six distinct linear solves and submitted five tree nodes. Lean checks that tree and proves the minimum-cost conclusion for every selection. This is not an experiment that visited sixteen million cases. It is one structured problem with a short exhaustive mathematical argument.

A second 24-candidate example puts every available correction on the first coordinate axis of $\Q^2$ and sets the residual to $(0,1)$. The row $(0,1)$ supplies a one-node certificate rejecting all selections. This illustrates compressed reasoning, not a claim that arbitrary difficult search problems have short certificates.

\section{Performance and semantic limits}

The earlier eight-candidate exhaustive producer remains available as a reference. The new producer accepts up to 24 candidates, with explicit limits on nodes and linear solves. These are implementation resource guards, not mathematical restrictions on the tree theorem.

In a favorable instance, separators or cost bounds eliminate large intervals. In the worst case the tree can still approach the full binary decision tree. Neither polynomial-time search nor a polynomial-size certificate is established. With an explicit admissibility table, checking an exclusion leaf also costs work proportional to that table.

The nonnegative cost requirement is essential to the implemented bound rule. If an optional candidate had negative cost, adding it could lower the cost below the lower selection's cost. The input schema rejects negative weights instead of silently applying an invalid monotonicity argument.

The theorem operates on the original rational matrices and declared subset policy. It does not prove that these finite data capture a native arithmetic, analytic, or representation-theoretic construction. That remains the semantic realization obligation of the domain application. In Komorebi Space, the admissible family must retain its declared meaning. A compressed certificate cannot replace that meaning with a more convenient problem.

\appendix
\chapter{Formal correspondence}\label{app:correspondence}
\section{How the prose matches the source}
The original formal source is the supplied \emph{Constructive Repair Lean v0.2.0} release~\cite{formal}. The namespace is \code{KomorebiRepair}. The notation $Q_D$, $q_D$, $O_D(r)$, and $E_S$ corresponds to \code{ObstructionSpace}, \code{obstructionMap}, \code{obstructionSpan}, and \code{menuMap}, respectively.

The mathematical statements here use $\exists u,w$ with the same sign convention as the source. Families quantify $\forall i\,\exists u_i,w_i$. They do not assume one shared correction vector. The algebraic base is a commutative ring. The dimension and section results specialize to a field, and finiteness is stated when numerical dimensions or terminating enumeration are used.

\begin{center}
\small
\begin{tabularx}{\linewidth}{@{}p{0.24\linewidth}Xp{0.16\linewidth}@{}}
\toprule
Mathematical topic & Principal formal module & Discussion\\
\midrule
Correction and repair & \code{Basic} & Chapter~\ref{ch:basic}\\
Free minimum & \code{Minimal} & Chapter~\ref{ch:minimal}\\
Legal subspace & \code{ConstrainedMinimal} & Chapter~\ref{ch:minimal}\\
Finite menu span & \code{Menu} & Chapter~\ref{ch:menu}\\
Invariant obstruction & \code{Conservation} & Chapter~\ref{ch:menu}\\
Faithful transport & \code{Transport} & Chapter~\ref{ch:transport}\\
Executable search & \code{Search} & Chapter~\ref{ch:search}\\
Quadratic digit step & \code{Arithmetic} & Chapter~\ref{ch:arithmetic}\\
Nonlinear boundaries & \code{Nonlinear} & Chapter~\ref{ch:limits}\\
Concrete witnesses & \code{Examples} & Chapters~\ref{ch:menu}, \ref{ch:transport}, \ref{ch:search}\\
\bottomrule
\end{tabularx}
\end{center}

A correspondence records the mathematical content of the formal declaration. It does not make every explanatory sentence or supplementary example a separately compiled theorem. In particular, \cref{cor:rank-difference}, the nonfield section counterexample, the span faithfulness illustration, and the enumeration bound are ordinary derivations given in the text. The source map distinguishes these from named Lean declarations.

\section{Constructive existence and executable construction}
The free minimum uses a linear section of a quotient. Its formal proof invokes a standard vector space splitting theorem and classical choice. The proof establishes an actual mathematical map with the required properties, but does not extract a canonical or efficiently computable representative system.

The finite optimizer is executable for the finite types in its contract. It returns a selected set, not a basis of the obstruction space or a coefficient table. Its logical correctness proves that suitable correction witnesses exist. The two evaluated finite examples are proved by kernel reduction, while their separate printed outputs serve only as execution traces.

These uses of the word ``constructive'' should not be merged. The first constructs existence within classical mathematics. The second defines and certifies a terminating finite algorithm. Chapters~\ref{ch:rational} and~\ref{ch:matrix-attainment} now supply a separate rational certificate checker and an explicit matrix construction. This is not a proof that the earlier Python optimizer and finite-domain Lean solver implement the same search algorithm.

\section{Stable source locators}
Appendix~\ref{app:register} gives every named theorem in the source registry, its source line, and its role in this monograph. The line numbers refer to the unchanged files in the accompanying original archive. Mathematical numbering in this document is independent of source declaration order.

The register was compared with the actual source files and with the theorem axiom entries in the supplied verification record. All 49 entries are represented. The declaration count includes helper facts and concrete examples. It is not a claim that the document establishes 49 independently novel mathematical discoveries.

\chapter{Original declaration register}\label{app:register}
Every source path below is relative to the root of the formal release. Prefix theorem names with \code{KomorebiRepair}. For the \code{Arithmetic}, \code{Nonlinear}, and \code{Examples} modules, retain that additional namespace component. The register is generated from the supplied JSON declaration list and linked to the corresponding mathematical discussion.

\moduleheading{Basic}{reg:Basic}

\record{obstructionMap_correction}{31}{Original corrections map to zero in the quotient. See \cref{thm:correctable}.}

\record{correctable_iff_obstruction_zero}{35}{Exact correction is equivalent to vanishing of the quotient class. See \cref{thm:correctable}.}

\record{repairable_iff_obstruction_mem_range}{46}{Added repair is equivalent to membership in the added quotient image. See \cref{thm:single}.}

\record{obstruction_mem_span}{65}{Each required residual class belongs to its declared obstruction span. See \cref{def:span}.}

\record{family_repairable_iff_span_le}{69}{All family equations are solvable exactly when the span is contained in the quotient image. See \cref{thm:family}.}

\record{reconstruct_cancellation}{82}{A quotient witness supplies exact cancellation in the original target. See \cref{cor:reconstruct}.}

\moduleheading{Minimal}{reg:Minimal}

\record{repair_dimension_lower_bound}{14}{Every finite dimensional added space has dimension at least that of the obstruction span. See \cref{thm:lower}.}

\record{exists_obstruction_section}{22}{The vector space quotient admits a classical linear section. See \cref{lem:section}.}

\record{exists_attaining_repair}{30}{Restricting a section yields a repair whose domain is the required span. See \cref{thm:attain}.}

\record{free_minimal_repair_rank}{46}{Attainment and the lower bound establish the least possible added dimension. See \cref{thm:minimum}.}

\moduleheading{ConstrainedMinimal}{reg:ConstrainedMinimal}

\record{admissible_subspace_attainment}{11}{A feasible legal subspace contains an attaining lift of the required span. See \cref{thm:legal}.}

\record{admissible_subspace_no_go}{30}{An excluded obstruction span rules out every added map valued in the legal subspace. See \cref{thm:legal-no-go}.}

\moduleheading{Menu}{reg:Menu}

\record{menuMap_apply}{18}{The menu map is the actual finite sum of selected coefficient directions. See \cref{eq:menu}.}

\record{menu_quotient_range}{24}{The quotient image of the menu is the span of the selected quotient directions. See \cref{lem:menu-range}.}

\record{finite_menu_repair_iff_span}{40}{Menu feasibility is equivalent to inclusion of the residual obstruction span. See \cref{thm:menu}.}

\record{menu_mono}{47}{Adding selected directions enlarges their quotient span. See \cref{prop:menu-mono}.}

\record{repairable_menu_mono}{53}{Feasibility is preserved when a menu is enlarged. See \cref{prop:menu-mono}.}

\moduleheading{Conservation}{reg:Conservation}

\record{invariant_preserved}{10}{Original and added corrections leave the stated linear invariant unchanged. See \cref{prop:invariant}.}

\record{conservation_no_go}{17}{A nonzero conserved residual invariant excludes a repair. See \cref{thm:conservation}.}

\record{menu_preserves_invariant}{28}{An invariant that kills every selected direction kills the whole menu map. See \cref{cor:menu-conservation}.}

\record{conservation_menu_no_go}{36}{The invariant obstruction applies to the actual selected coefficient map. See \cref{cor:menu-conservation}.}

\moduleheading{Transport}{reg:Transport}

\record{correction_range_maps}{12}{A commuting correction square transports the original image into the target image. See \cref{lem:range-maps}.}

\record{transport_obstruction_naturality}{23}{The induced quotient map carries the actual residual class to its transported class. See \cref{prop:naturality}.}

\record{faithfulOn_iff_injective}{32}{Faithfulness on a submodule is injectivity of the restricted quotient map. See \cref{prop:faithful-inj}.}

\record{transport_faithfulness_on_obstruction_span}{49}{Faithfulness on the required span makes target correctability reflect source correctability. See \cref{thm:faithful}.}

\record{nonfaithful_witness}{63}{Failure of faithfulness produces a specific nonzero class that disappears. See \cref{prop:loss-witness}.}

\moduleheading{Search}{reg:Search}

\record{bestByCheck_spec}{20}{The filtered finite selector returns a passing candidate of minimum cost. See \cref{lem:best}.}

\record{bestByCheck_none_iff}{30}{The selector rejects exactly when no listed candidate passes its Boolean predicate. See \cref{lem:best}.}

\record{repairMenuCheck_spec}{45}{The executable Boolean is equivalent to the original family repair equations. See \cref{prop:check}.}

\record{mem_allMenus}{54}{Every subset of the finite candidate universe occurs in the enumerated list. See \cref{lem:allmenus}.}

\record{menu_optimality_certificate}{70}{A selected menu is allowed, feasible, and minimum cost in the declared universe. See \cref{thm:optimizer}.}

\record{menu_search_none_iff}{86}{Rejection is complete within the finite candidate and admissibility contract. See \cref{thm:none}.}

\moduleheading{Arithmetic}{reg:Arithmetic}

\record{Arithmetic.linear_digit_exists}{13}{Bezout data construct an integer solution of the new linear digit congruence. See \cref{lem:digit}.}

\record{Arithmetic.quadratic_digit_correction}{20}{The exact quadratic expansion proves one higher divisibility scale. See \cref{lem:quad}.}

\record{Arithmetic.quadratic_hensel_step}{33}{A coprime derivative supplies a next quadratic correction at a prime power scale. See \cref{thm:hensel}.}

\record{Arithmetic.hensel_seven_trace}{43}{The seven recorded representatives satisfy their seven exact integer congruences. See \cref{sec:trace}.}

\record{Arithmetic.singular_mod_four_impossible}{52}{There is no square root of two in the ring of integers modulo four. See \cref{prop:mod4}.}

\record{Arithmetic.singular_derivative_can_still_lift}{54}{The exact root one of the equation $x^2=1$ persists at every dyadic depth. See \cref{prop:singular-success}.}

\moduleheading{Nonlinear}{reg:Nonlinear}

\record{Nonlinear.nonlinear_root_despite_first_order_obstruction}{10}{The derivative obstruction at zero does not exclude the displayed nonlinear root. See \cref{prop:nonlinear-root}.}

\record{Nonlinear.linearized_solution_without_real_root}{24}{An actual solvable derivative equation does not imply a global real root. See \cref{prop:no-real-root}.}

\moduleheading{Examples}{reg:Examples}

\record{Examples.first_menu_entry_insufficient}{24}{The first rational menu direction alone cannot cancel the residual. See \cref{ex:menu}.}

\record{Examples.second_menu_entry_insufficient}{33}{The second rational menu direction alone cannot cancel the residual. See \cref{ex:menu}.}

\record{Examples.two_menu_entries_repair}{39}{Two selected rational directions supply the displayed exact correction. See \cref{eq:two-entry-witness}.}

\record{Examples.one_free_direction_repairs}{44}{One freely chosen rational direction repairs the same residual. See \cref{ex:menu}.}

\record{Examples.coarse_information_loss}{50}{Projection sends an uncorrectable rational residual to a correctable target residual. See \cref{ex:loss}.}

\record{Examples.collapsing_transport_not_faithful}{59}{The associated quotient transport fails the required span faithfulness condition. See \cref{ex:loss}.}

\record{Examples.executable_menu_returns_both}{81}{Kernel reduction proves the optimizer returns both finite field directions. See \cref{sec:f2}.}

\record{Examples.executable_restricted_menu_rejects}{88}{Kernel reduction proves rejection under the one entry admissibility rule. See \cref{sec:f2}.}

\record{Examples.executable_menu_is_globally_optimal}{96}{The generic optimum theorem proves the two entry lower bound for every feasible subset. See \cref{sec:f2}.}

\chapter{Verification and provenance}\label{app:evidence}
\section{The evidence attached to the formal release}
The v0.2.0 verification record reports successful compilation of ten proof modules, an axiom audit of all 49 named theorem declarations, and individual replay of those modules through Lean's kernel. The allowed transitive axiom set is
\[
\{\code{propext},\ \code{Classical.choice},\ \code{Quot.sound}\}.
\]
No additional assumption axiom, placeholder axiom, or native decision axiom is allowed in the accepted theorem library.

The environment recorded by the release is Lean \code{4.34.0-rc2}. Its commit is
\begin{center}\footnotesize\nolinkurl{6a10ac8c22beadecabdbb0919c2b50214762f91d}\end{center}
with Mathlib revision
\begin{center}\footnotesize\nolinkurl{aeb72f1e9981b6d35b426b16dca7759be3e7ba06}.\end{center}
All dependency pins remain in the original \code{lake-manifest.json}.

The recorded \code{leanchecker} replay uses Lean's own kernel against its imported environment. It is not an independently implemented checker and is not a source rebuild of every imported Mathlib theorem. The small source archive does not bundle the compiler or its full dependency cache.

\begin{center}
\small
\begin{tabularx}{\linewidth}{@{}Xr@{}}
\toprule
Recorded formal release check & Coverage\\
\midrule
Proof module compilations & 10\\
Named theorem axiom audits & 49\\
Individual same kernel module replays & 10\\
Intended compiler rejection controls & 5\\
Separate forbidden-axiom policy control & 1\\
Python unit tests in the recorded replay & 20\\
Individual cases in the recorded Python audit & 3,307\\
Explicit assertions in that audit & 71\\
\bottomrule
\end{tabularx}
\end{center}

The original archive's final extraction check also records successful rebuilding and replay after extraction. These are historical execution records carried with the unchanged source. Those historical receipts remain unchanged. The expanded v0.3.0 execution has its own fresh compilation and replay record, described in the additional verification appendix.

\section{Negative controls and their meaning}
The five compiler controls intentionally assert a false equality, omit required faithfulness, claim an insufficient menu succeeds, claim a false divisibility gain, or introduce a proof placeholder. The verifier requires the intended diagnostics, rather than accepting an unrelated parsing or missing instance failure as a successful negative test.

A sixth control deliberately adds a false axiom. Lean can correctly verify a theorem conditional on that axiom. The separate release policy then rejects the dependency. This distinguishes logical checking under assumptions from acceptance of those assumptions as part of a mathematical release.

The development evidence retains failed proof drafts and superseded verification runs. In particular, the release ledger records corrections to an axiom-name expectation and to a negative control that initially failed for a missing decidability instance. Only the hardened final record governs the stated formal status.

\section{Experimental evidence is a different layer}
The earlier exact audit~\cite{audit} includes rational and finite-field linear calculations, Euler inverse identities, selected cyclotomic coefficients, finite transport graphs, and arithmetic and nonlinear controls. Its 3,307 case records are retained in the formal handoff. These computations do not prove the unrestricted versions of the surrounding arithmetic or representation-theoretic problems.

The finite graph and Heisenberg matrix calculations are not a formal construction of a local-field representation. The Euler and cyclotomic tests are not a new proof of the full Euler conjecture. The original ten Lean modules do not contain an integrated replay of Papers III or XI. No result here resolves Paper XI's Conjecture 66 or provides a universal arithmetic repair theorem.

\section{The monograph edition}
The source of this expanded edition is a standalone LaTeX file. It contains the bibliography and all commutative diagrams, with no external image or font files required. The original formal release is included unchanged alongside the document.

The unchanged v0.2.0 formal source archive has SHA-256
\begin{center}\footnotesize\nolinkurl{d906e98ddc2a661f3201ba51e1f1049a49dfa9326078384dfa851b608b80f337}.\end{center}
Its checksum agrees with the recorded extraction report. The twelve positive source and audit-harness files listed in the verification record match their stored SHA-256 values. The declaration registry matches the theorem names in those source files and their recorded axiom entries.

The publication evidence separates these integrity checks from LaTeX compilation and layout inspection. It records the resulting page count and document hashes. The ten original proof modules were not changed. New matrix modules and fresh proof evidence are recorded separately. No mathematical completion claim is inferred merely from producing a PDF.

\chapter{Rational bridge declaration register}\label{app:rational-register}

The twenty general declarations below supplement the unchanged original register. All source paths are relative to the release root. The full machine-readable register in \code{docs/theorem_registry.json} includes every concrete case and every named rejection theorem as well.

\section{RationalCertificates}

Source: \nolinkurl{lean/KomorebiRepair/RationalCertificates.lean}.

\record{RationalCertificates.checkRepair_sound}{45}{An accepted rational coefficient table repairs the fixed original residual family.}

\record{RationalCertificates.annihilator_kills_residual}{51}{An annihilator of both correction ranges annihilates every successfully repaired residual.}

\record{RationalCertificates.separator_sound}{64}{A nonzero detected residual rules out a correction by that fixed menu.}

\record{RationalCertificates.checkSeparator_sound}{69}{The Boolean separating-row checker is sound.}

\record{RationalCertificates.checkOptimal_sound}{79}{A successful menu with separators for all cheaper admissible menus is globally minimum cost within the declared menu.}

\record{RationalCertificates.checkNoRepair_sound}{96}{Separators for every admissible subset certify rejection in that finite universe.}

\record{RationalCertificates.dimension_certificate_right_inverse}{104}{Every competing repair supplies an explicit right inverse of the detected added map.}

\record{RationalCertificates.width_le_of_right_inverse}{121}{A rectangular rational map with a right inverse has at least as many source as target coordinates.}

\record{RationalCertificates.dimension_certificate_lower_bound}{129}{The matrix identities imply the lower bound against every competing rational augmentation.}

\record{RationalCertificates.checkDimension_sound}{142}{A successful repair and a lower certificate certify attained minimum dimension.}

\record{RationalCertificates.checkRank_sound}{157}{A factorization and an identity minor certificate prove the declared exact rank.}

\section{ConstructiveMatrix}

Source: \nolinkurl{lean/KomorebiRepair/ConstructiveMatrix.lean}.

\record{RationalCertificates.split_constructs_repair}{13}{The split and residual factorization explicitly construct the correction coefficients.}

\record{RationalCertificates.split_correctable_iff}{30}{The complete detector vanishes exactly when the original rational matrix equation is solvable.}

\record{RationalCertificates.checkOptimalConstruction_sound}{58}{Four matrix identities construct an attaining repair and prove its minimum added dimension.}

\record{RationalCertificates.constructed_repair_in_legal_image}{77}{An additional image identity certifies that the constructed directions are legal.}

\record{RationalCertificates.decodeCell_zero}{89}{A typed rational cell with denominator zero is rejected.}

\record{RationalCertificates.decodeCell_nonzero}{92}{A cell with nonzero denominator decodes to the stated rational value.}

\record{RationalCertificates.decodeCell_sound}{96}{Successful cell decoding entails a nonzero denominator and the correct value.}

\section{Wire}

Source: \nolinkurl{lean/KomorebiRepair/Wire.lean}.

\record{RationalCertificates.checkWireMatrix_sound}{20}{The finite matrix validity check certifies all denominators.}

\record{RationalCertificates.decodeWireMatrix_cell}{23}{Matrix decoding agrees pointwise with successful cell decoding.}

\section{Concrete exported problems}

Each row below denotes one original problem and its separately supplied certificate. Every case has six named proofs: \code{input_encoding_valid}, \code{certificate_encoding_valid}, \code{certificate_checked}, \code{original_equations_checked}, \code{output_matches_construction}, and \code{optimal}. The last theorem quantifies over every competing rational augmentation, not only the candidates tested by Python.

\begin{longtable}{@{}lrrrrrp{1.75in}@{}}

\toprule Case & $m$ & $n$ & $s$ & $d$ & $t$ & Example class\\ \midrule\endhead

Case000 & 3 & 1 & 2 & 2 & 1 & worked\\

Case001 & 2 & 1 & 1 & 1 & 1 & fractional\\

Case002 & 0 & 2 & 3 & 0 & 0 & empty target\\

Case003 & 3 & 0 & 2 & 3 & 2 & empty original\\

Case004 & 3 & 1 & 0 & 2 & 0 & empty residuals\\

Case005 & 3 & 1 & 2 & 2 & 0 & zero residual\\

Case006 & 3 & 3 & 2 & 0 & 0 & surjective\\

Case007 & 3 & 2 & 2 & 3 & 2 & zero original matrix\\

Case008 & 0 & 0 & 0 & 0 & 0 & all empty\\

Case009 & 4 & 2 & 2 & 3 & 2 & nontrivial rank two\\

Case010 & 3 & 1 & 2 & 3 & 2 & random 000\\

Case011 & 3 & 3 & 1 & 0 & 0 & random 005\\

Case012 & 3 & 2 & 2 & 2 & 2 & random 006\\

Case013 & 3 & 4 & 3 & 3 & 3 & random 007\\

Case014 & 2 & 2 & 0 & 0 & 0 & random 008\\

Case015 & 2 & 1 & 2 & 1 & 1 & random 010\\

Case016 & 3 & 2 & 0 & 3 & 0 & random 013\\

Case017 & 0 & 3 & 4 & 0 & 0 & random 014\\

Case018 & 2 & 3 & 4 & 0 & 0 & random 015\\

Case019 & 1 & 4 & 1 & 1 & 0 & random 016\\

Case020 & 3 & 1 & 3 & 2 & 2 & random 017\\

Case021 & 4 & 4 & 0 & 0 & 0 & random 018\\

Case022 & 0 & 1 & 1 & 0 & 0 & random 019\\

Case023 & 1 & 4 & 4 & 1 & 1 & random 020\\

\bottomrule\end{longtable}

The input and certificate sources are grouped in \nolinkurl{Generated/Inputs00.lean} through \nolinkurl{Inputs03.lean}, and \nolinkurl{Generated/Certificates00.lean} through \nolinkurl{Certificates03.lean}. Groups contain six cases each. The JSON index gives the exact original input hash for each row.

\section{Menu and corruption examples}

The eight \code{RationalMenus} declarations check the two-entry optimum, restricted rejection, invalid data, and zero- and one-rank certificates. Six further named proofs in \code{RejectedCertificates.lean} establish that deliberately corrupted checks return false. Their companion invalid files attempt to use these false equalities as true equalities and are required to fail for the intended type mismatch.

\chapter{Rational bridge verification and use}\label{app:fresh}

\section{The v0.3.0 checkpoint}

The following records describe the preceding release. Its source modules are retained unchanged and freshly checked again in the v0.4.0 run. The v0.4.0 menu extension is documented in the next appendix. A later appendix records the v0.5.0 tree extension, followed by the present release contract.

The runtime was recovered from the previously supplied archive parts and its joined checksum was matched before extraction. The compiler reports Lean \code{4.34.0-rc2}; the actual Mathlib checkout has revision \nolinkurl{aeb72f1e9981b6d35b426b16dca7759be3e7ba06}. The combined project was rebuilt from its own source files after removing its previous object files. The dependency cache was retained.

\begin{center}\begin{tabularx}{\linewidth}{@{}Xr@{}}\toprule Fresh check & Result\\\midrule

Proof modules compiled and individually replayed & 23 / 23\\

Named theorem dependencies audited & 227 / 227\\

Unchanged original declarations & 49\\

New general certificate and decoding declarations & 20\\

New concrete declarations & 158\\

Exported original problems with optimality proofs & 24 / 24\\

Paired corruption rejection controls & 6 / 6\\

Separate forbidden-axiom policy control & 1 / 1\\

Lean input echoes compared with original JSON & 24 / 24\\

New Python unit tests & 34 / 34\\

Exact matrix cases with independent SymPy ranks & 264 / 264\\

Inherited Python unit tests rerun & 20 / 20\\

Inherited exact audit case records reproduced & 3,307\\

\bottomrule\end{tabularx}\end{center}

The 158 new concrete declarations consist of 144 input/certificate proofs for the 24 original problems, eight menu/rank examples, and six rejection examples. These are coverage counts, not counts of independent discoveries. The general theorems and the concrete proofs have separately labeled entries in the source registry.

Every audited dependency set is contained in \{\code{propext}, \code{Classical.choice}, \code{Quot.sound}\}. No added assumption, placeholder or native-decision axiom is accepted. The module replay invokes the same Lean kernel, not an independently implemented proof checker. It does not re-prove all imported Mathlib declarations.

\section{The single-command example}

The release includes \code{scripts/certify_one.py}. Given an original JSON input and the pinned offline kit, it constructs a certificate in a new local directory, emits the input separately from the witness, compiles the generic modules and the concrete proof, audits all relevant axioms, replays the modules, and compares a Lean-decoded echo with the original numeric data.

\begin{Verbatim}[fontsize=\small,breaklines=true]

export KOMOREBI_LEAN_KIT=/path/to/Komorebi_Portable_Lean_4.34.0-rc2

python scripts/certify_one.py examples/corpus/Case001_input.json --output /tmp/my-repair

\end{Verbatim}

This exact workflow was executed on the fractional example. It ended with \nolinkurl{KERNEL_CHECKED_OPTIMAL_REPAIR} and $t=1$. It requires no network access or publication service. An existing output directory is refused so that previous evidence is not silently overwritten. The large compiler and dependency kit are not duplicated in the small source handoff.

\section{Failures and their corrections}

The first compilation attempt with a broad tactic import was interrupted. The implemented modules use the required direct imports. Plain \code{decide} did not reduce the first rational matrix fixture, so the producer now emits \code{decide +kernel}. The typed validity predicate initially lacked a reducible decidability path and was corrected. An ambiguous corruption diagnostic was replaced by paired rejection proofs. In the final control, the false checker result is itself proved before an invalid true result is rejected.

The development snapshots and logs remain in \code{evidence/development}. Interrupted and superseded runs are not counted as successful checks. The inherited v0.3.0 record is preserved under \nolinkurl{evidence/verification/v0.3.0_inherited}. The v0.4.0 receipt is preserved under \nolinkurl{evidence/verification/v0.4.0_inherited}. That historical combined receipt is now preserved under \nolinkurl{evidence/verification/v0.5.0_inherited}. The current receipt governs the release candidate described in the final appendix.

\section{Remaining boundaries}

The new code is a standalone extension. It has not been integrated into the authoritative Komorebi Python library or Paper III source tree. At that checkpoint, the rational-menu exporter was still pending. Chapter 10 now supplies that path. Full formalization of the external JSON parser and a formal connection of the matrix rank-difference formula to the earlier quotient API remain separate tasks. The current free-augmentation exporter, certificate soundness, attaining construction, and finite literal proofs are complete within their stated rational matrix scope.

The construction does not establish unrestricted nonlinear repair, integral or arithmetic realization, or a universal semantic theorem for Komorebi Space. It does not alter Conjecture 66. Its purpose is narrower and operational: allow an untrusted numerical search to submit exact witnesses that are checked against the original problem.

\chapter{Menu certificates: sources and fresh verification}
\label{app:menu-verification}

This appendix preserves the v0.4.0 continuation record, rather than describing the full v0.5.0 count. The 23 proof modules from v0.3.0 remain byte-for-byte unchanged. Fourteen additional modules supply a seven-theorem cover interface, 24 exported rational-menu problems, and eight paired corruption controls. The full combined project is compiled and replayed again, rather than accepting the prior receipts as a substitute for the new run.

\section{General declarations}

The new general source is \nolinkurl{lean/KomorebiRepair/RationalMenuCover.lean}. The following records give its exact current line locations. Definitions, decidability instances, and diagnostic JSON display functions are not counted as theorems.

\record{KomorebiRepair.RationalCertificates.mul_masked_entry}{16}{A column-mask identity for exact matrix multiplication.}
\record{KomorebiRepair.RationalCertificates.separates_subset}{22}{A valid separator for a larger menu separates every subset.}
\record{KomorebiRepair.RationalCertificates.covered_infeasible}{48}{A covered selection is infeasible when every listed block is valid.}
\record{KomorebiRepair.RationalCertificates.checkCoveredOptimal_sound}{59}{Acceptance proves admissibility, the original equations, and minimum cost among admissible feasible selections.}
\record{KomorebiRepair.RationalCertificates.checkCoveredNoRepair_sound}{77}{Acceptance excludes every selection admitted by the original rule.}
\record{KomorebiRepair.RationalCertificates.tableAllowed_iff}{94}{The Boolean admissibility table is equivalent to membership in its fixed input list.}
\record{KomorebiRepair.RationalCertificates.additiveCost_empty}{97}{The empty selection has zero additive cost.}

\section{Exported original problems}

Every input is emitted independently of its coefficient and separation witnesses. The six input modules and six certificate modules under \nolinkurl{MenuGenerated/} hold four cases each. The full source registry, JSON files, and hashes are included in the package.

\begin{longtable}{@{}p{0.12\linewidth}p{0.44\linewidth}p{0.15\linewidth}rr@{}}
\toprule Case & Main test & Outcome & Cost & Blocks\\\midrule
\endfirsthead
\toprule Case & Main test & Outcome & Cost & Blocks\\\midrule
\endhead
Menu000 & two entry optimum & optimum & 2 & 2\\
Menu001 & size one restriction & none & -- & 2\\
Menu002 & pair cheaper than singleton & optimum & 2 & 2\\
Menu003 & fractional witness & optimum & 7 & 1\\
Menu004 & full infeasible cover & none & -- & 1\\
Menu005 & free redundant directions & optimum & 3 & 1\\
Menu006 & zero cost tie & optimum & 0 & 0\\
Menu007 & nonhereditary allowed & optimum & 2 & 0\\
Menu008 & empty admissible universe & none & -- & 0\\
Menu009 & empty residual family & optimum & 0 & 0\\
Menu010 & zero target dimension & optimum & 0 & 0\\
Menu011 & zero original domain & optimum & 2 & 2\\
Menu012 & zero candidates correctable & optimum & 0 & 0\\
Menu013 & zero candidates blocked & none & -- & 1\\
Menu014 & multiple residuals & optimum & 6 & 2\\
Menu015 & inadmissible free solution & optimum & 1 & 0\\
Menu016 & seeded 000 & optimum & 3 & 0\\
Menu017 & seeded 001 & optimum & 0 & 0\\
Menu018 & seeded 002 & optimum & 1 & 1\\
Menu019 & seeded 003 & optimum & 0 & 0\\
Menu020 & seeded 004 & optimum & 0 & 0\\
Menu021 & seeded 005 & optimum & 0 & 0\\
Menu022 & seeded 006 & optimum & 4 & 2\\
Menu023 & seeded 007 & optimum & 1 & 0\\
\bottomrule\end{longtable}

There are twenty positive optimal-menu results and four no-admissible-menu results. One of the latter has an empty admissibility table. Its rejection is explicitly vacuous, not a claim that a nonzero cokernel class has been detected.

\section{Fresh verification results}

\begin{center}\begin{tabularx}{\linewidth}{@{}Xr@{}}\toprule Check & Result\\\midrule
Combined proof modules compiled and replayed & 37 / 37\\
Named theorem dependencies audited & 378 / 378\\
Unchanged v0.3.0 declarations & 227\\
Additional general declarations & 7\\
Additional concrete declarations & 144\\
New original menu problems exported to Lean & 24 / 24\\
New matrix, admissibility, and cost echoes & 24 / 24\\
New paired corruption controls & 8 / 8\\
Previous paired corruption controls rerun & 6 / 6\\
Separate forbidden-axiom policy control & 1 / 1\\
New Python unit tests & 50 / 50\\
Exact menu problems independently compared & 120 / 120\\
Individual subset rank comparisons & 709 / 709\\
Previous rational Python unit tests rerun & 34 / 34\\
Original Python unit tests rerun & 20 / 20\\
Original exact audit records reproduced & 3,307\\
\bottomrule\end{tabularx}\end{center}

The 144 additional concrete declarations comprise 136 input and certificate statements for the 24 problems and eight positive proofs that corrupted checks return false. This is a coverage count, not a count of new mathematical discoveries. The underlying finite linear alternative, monotonicity of linear span, and finite optimization arguments are classical.

All transitive axiom lists are contained in \{\code{propext}, \code{Classical.choice}, \code{Quot.sound}\}. Concrete certificates use \code{decide +kernel}. Replay invokes Lean's own kernel against the pinned imports. It is not an independently implemented checker and does not rebuild or re-prove all Mathlib dependencies.

The independent Python comparison uses SymPy matrix ranks to assess every subset of each recorded problem. It checks the producer's selected minimum or infeasibility outcome, rather than merely comparing two calls to the same elimination routine. The seeded run contains 120 problems and 709 individual subset comparisons. All rational arithmetic is exact.

\section{Rejection controls and command execution}

The new corruption controls remove cheaper-menu coverage, replace a separator by zero, overstate a block's scope, change correction coefficients, propose an expensive singleton despite a cheaper repair, select a forbidden menu, claim rejection of a feasible family, and attempt to use an unselected column. Each corrupted Boolean check is first proved false. A separate invalid source then tries to use that result as true and must fail with the intended type mismatch. An unrelated missing instance or syntax failure is not accepted as a rejection result.

The single-command workflow was executed on the fractional problem and on the four-candidate globally infeasible menu. The former returns \nolinkurl{KERNEL_CHECKED_OPTIMAL_MENU}, cost seven, with coefficient $-5/4$. The latter returns \nolinkurl{KERNEL_CHECKED_NO_ADMISSIBLE_MENU}, using one separator for sixteen subsets. Both commands compile the sources, audit all relevant theorem dependencies, replay the modules, and match the Lean input echo to the original matrices and constraints.

A separate trial-budget control deliberately permits zero trials. It returns \code{UNRESOLVED} before claiming a mathematical outcome. The log is retained. This distinction is part of the operational contract, not merely a closing qualification.

\section{Development record and remaining scope}

The first new cover-module compilation failed because \code{if_true} and \code{if_false} are deprecated in the pinned runtime and warnings are treated as errors. They were replaced by \code{ite_true} and \code{ite_false}; the subsequent source compiled. That failed attempt and its exact source snapshot remain in the development evidence.

The existing free-repair command was also adjusted to select only the general declarations it actually imports, rather than every declaration later assigned the general-layer label. Its mathematical modules are unchanged. This prevents an expanded registry from causing the older command to request nonexistent imports.

The external JSON parser, source emission, hashing, and file system are tested but not formally verified. The checked mathematical statements apply to the literal typed input compiled by Lean. Separate input emission, hashes, and a Lean-decoded echo provide executed evidence that those literals match the intended problem. This boundary remains explicit.

The v0.4.0 producer was exhaustive and bounded. Its implementation remains unchanged as a reference; the following continuation supplies a separate branch-and-bound path. It does not minimize evidence size. No integration into the authoritative Komorebi Python package or Paper III source tree is claimed. No theorem here constructs native arithmetic hosts, proves nonlinear convergence, or establishes a universal realization theorem for Komorebi Space. Those tasks remain separate from the now-operational costed rational-menu certificate path.


\chapter{Branch certificate verification}

\label{app:branch}

This appendix records the v0.5.0 continuation. The full 37-module v0.4.0 proof tree is retained byte-for-byte. Fourteen additional modules supply the general branch checker, 24 exported problems, and paired corruption controls. Earlier appendices remain historical records of their respective releases.

\section{General declaration correspondence}

Source: \nolinkurl{lean/KomorebiRepair/BranchCertificate.lean}. Each declaration is proved with explicit matrix, cost and admissibility data. No field assumes that an optimal repair exists.

\record{KomorebiRepair.RationalCertificates.intervalExcluded_sound}{42}{A checked exclusion of an interval from the original optional admissibility table rules out every admitted selection in that interval.}

\record{KomorebiRepair.RationalCertificates.additiveCost_mono}{53}{Nonnegative additive column costs are monotone under inclusion of the selected set.}

\record{KomorebiRepair.RationalCertificates.boundCheck_mono}{68}{A checked lower cost bound propagates to a larger cost. No threshold means impossibility, so a cost leaf cannot apply.}

\record{KomorebiRepair.RationalCertificates.checkBranch_sound}{95}{Structural induction on an accepted tree gives a lower cost bound or infeasibility for every admitted feasible selection in its interval.}

\record{KomorebiRepair.RationalCertificates.checkTreeOptimal_sound}{131}{A successful original coefficient equation and an accepted root tree certify minimum cost within the declared candidate universe.}

\record{KomorebiRepair.RationalCertificates.checkTreeNoRepair_sound}{148}{An accepted root tree in impossibility mode excludes every selection admitted by the original rule.}

\section{Exported case register}

The following cases are emitted from independent original JSON inputs and exact witness records. The tree-node count is the submitted certificate size, not a count of separately evaluated subsets. Cost is omitted for a rejection result.

\begin{longtable}{@{}p{0.14\linewidth}p{0.43\linewidth}rrr@{}}

\toprule Case & Description & Candidates & Cost & Nodes\\\midrule\endhead

Branch000 & two entry optimum & 2 & 2 & 5\\

Branch001 & size one restriction & 2 & -- & 5\\

Branch002 & cheaper pair & 3 & 2 & 13\\

Branch003 & fractional & 1 & 7 & 3\\

Branch004 & zero cost & 2 & 0 & 5\\

Branch005 & nonhereditary only pair & 2 & 2 & 5\\

Branch006 & empty admissible table & 2 & -- & 1\\

Branch007 & empty residuals & 2 & 0 & 1\\

Branch008 & zero target & 2 & 0 & 1\\

Branch009 & zero domain & 2 & 2 & 5\\

Branch010 & no candidates feasible & 0 & 0 & 1\\

Branch011 & no candidates impossible & 0 & -- & 1\\

Branch012 & multiple residuals & 2 & 6 & 5\\

Branch013 & free solution forbidden & 2 & 1 & 5\\

Branch014 & duplicate costs & 3 & 3 & 7\\

Branch015 & incumbent improvement & 3 & 4 & 9\\

Branch016 & non downward rules & 3 & 5 & 11\\

Branch017 & four menu impossible & 4 & -- & 1\\

Branch018 & large 24 optimum & 24 & 5 & 5\\

Branch019 & large 24 impossible & 24 & -- & 1\\

Branch020 & large 24 zero cost & 24 & 0 & 3\\

Branch021 & large 24 finite admissibility & 24 & 5 & 5\\

Branch022 & large 24 empty table & 24 & -- & 1\\

Branch023 & nonempty rules empty residuals & 2 & 0 & 3\\

\bottomrule\end{longtable}

The register contains 18 optimal-menu certificates and 6 no-admissible-menu certificates. An empty admissibility table is explicitly vacuous. It is not evidence that a nonzero arithmetic obstruction was discovered.

\section{Fresh combined verification}

The current evidence record is \nolinkurl{evidence/verification/current/verification.json}. The prior current record is preserved under \nolinkurl{evidence/verification/v0.4.0_inherited}. The full combined library is recompiled from source and each module is replayed through Lean\textquotesingle s kernel against the pinned imported environment.

\begin{center}\begin{tabularx}{\linewidth}{@{}Xr@{}}\toprule Check & Result\\\midrule

Combined proof modules compiled and replayed & 51 / 51\\

Named theorem dependencies audited & 526 / 526\\

Unchanged v0.4.0 proof modules & 37 / 37\\

Additional general declarations & 6\\

Additional concrete declarations & 142\\

New exported original-input echoes & 24 / 24\\

New paired corruption controls & 10 / 10\\

Previous paired corruption controls rerun & 14 / 14\\

Separate forbidden-axiom policy control & 1 / 1\\

New Python unit tests & 50 / 50\\

New exact branch problems rechecked & 204 / 204\\

Small problems compared with independent exact ranks & 199 / 199\\

Individual subset rank comparisons & 3,173 / 3,173\\

Comparisons with the previous exhaustive solver & 199 / 199\\

Previous rational and menu unit tests rerun & 84 / 84\\

Original Python unit tests rerun & 20 / 20\\

Original audit case records reproduced & 3,307\\

\bottomrule\end{tabularx}\end{center}

The additional 142 concrete declarations are input-encoding, certificate, equation, cost and rejection proofs. They are not a count of independent mathematical discoveries. The six general results specialize classical finite optimization and linear obstruction arguments to the declared proof-tree interface.

All audited theorem dependencies are contained in \{\code{propext}, \code{Classical.choice}, \code{Quot.sound}\}. Concrete reductions use \code{decide +kernel}. No new assumption axiom, \code{sorryAx}, or native-decision axiom is admitted. Replay uses Lean\textquotesingle s own kernel, not a second independently implemented checker and not a fresh proof of all Mathlib dependencies.

The external comparison enumerates subsets only in the small benchmark oracle, never in the new production checker. SymPy supplies independent exact ranks for 3,173 subsets of 199 small problems. The earlier exhaustive producer agrees on the outcome and minimum cost. Ties need not return the same selected labels. Five larger structured problems are not exhaustively tested by that oracle; their exported trees receive the formal check instead.

\section{Controls, command execution, and provenance}

The new negative controls forge a root cost bound, falsely exclude an all-admitted interval, use a zero separator, repeat a resolved split, duplicate a cost child, swap branch meanings, alter correction coefficients, use an unselected column, choose a forbidden selection, and claim no repair despite a valid incumbent. Each corrupted checker result is proved false. A separate invalid source must then fail specifically when treating that equality as true. Missing instances and syntax errors do not count as successful mathematical rejection.

The large optimal and infeasible cases were executed through the single-command interface. The command independently emits the original matrices, weights and optional admissibility table, then the witnesses. Compilation, axiom auditing, same-kernel replay, and an exact Lean input echo follow. Unlike the old menu echo, the new echo does not enumerate every derived selection: it returns the exact original policy data.

A deliberately zero node budget produces \code{UNRESOLVED}, not a rejection of the mathematical problem. Rejected external evidence has the same unresolved status. A no-repair status requires an accepted certificate in the mode that forbids cost leaves.

The first branch-module proof used an incorrect Boolean conjunction lemma name. The pinned compiler rejected it. The corrected proof uses \code{Bool.and_eq_true_iff}; both source and log versions are preserved.

A later audit review found that inherited compilation labels used only file stems. Generated modules in distinct namespaces could therefore overwrite one another\textquotesingle s diagnostic log and check key. That run is preserved as superseded. The verifier now uses full relative module names, rejects duplicate command labels, and repeats the full combined verification. No theorem statement, proof assumption, or certificate was weakened by this correction.

A separate evaluator inspection caught a more consequential implementation issue. The first version wrote a bounded-looking proposition over the admissibility list and asked Lean to infer its decision procedure. Lean instead selected a finite-type universal quantifier over candidate subsets. The theorem remained sound, but an exclusion leaf could still perform a powerset scan. The implementation now calls \code{List.all} explicitly and proves its connection to the mathematical exclusion predicate. A 24-candidate empty-table regression and a nonhereditary 24-candidate table regression exercise this path. The final verifier also inspects the elaborated definition and rejects the previous hidden powerset quantifier. Both evaluator outputs and the superseded source are retained.

\section{Trust and remaining scope}

Typed rational decoding and the mathematical certificate theorems are formalized. The complete JSON parser, untrusted search, source emitter, hashing and filesystem are tested, not formally verified. The kernel certifies the literal typed input in the generated source. Separate source emission, hashes, and a complete original-data echo test its binding to the supplied JSON.

The producer accepts up to 24 candidates and has finite resource guardrails. Search and certificate size can remain exponential. The greedy order is not proved optimal and proof-tree size is not minimized. The certificate soundness theorem does not require the producer to exhaust its entire budget or follow a particular search order.

Direct integration into the authoritative Komorebi Python library and Paper III source tree, full parser verification, and a formal quotient-to-matrix rank-difference bridge remain separate tasks. There is no new native arithmetic host construction, nonlinear convergence theorem, change to Conjecture 66, or unrestricted Komorebi Space realization claim.

\chapter{Release contract and reproducibility}
\label{app:release-candidate}
\label{app:rc2-portability}

\section{A frozen mathematical component}

Release candidate 1.0.0rc3 is the standalone exact rational repair component described in this monograph. It does not assign a final version to all of Operation Komorebi. The 51 proof modules and 526 named theorem declarations are unchanged from v0.5.0. No new mathematical assertion is introduced by the release number.

The component constructs freely chosen augmentations, chooses minimum cost repairs from prescribed menus, and verifies branch certificates accounting for skipped alternatives. Every acceptance concerns the original matrices, coefficient ring, costs, and admissibility rules. No claim about integral coefficients, nonlinear convergence, native arithmetic realization, or Conjecture 66 follows merely from a successful rational matrix check.

The previous candidates' API corrections, encoding regression, and runtime pin checks are retained. Their detailed histories remain in the source archive's documents and logs. This appendix consolidates the current release contract rather than accumulating another chronological release report inside the mathematical exposition.

\section{Using the installed interface}

The Python wheel has no third party runtime dependencies. Its public interface accepts exact matrices and returns finite certificate data. For example:

\begin{Verbatim}[fontsize=\small,breaklines=true]
from komorebi_repair import ExactMatrix, menu_problem, solve, verify

D = ExactMatrix.of([[1], [0], [0]])
E = ExactMatrix.of([[0, 0], [1, 0], [1, 1]])
R = ExactMatrix.of([[0], [1], [0]])
p = menu_problem(D, E, R, costs=[1, 1])
answer = solve(p)
assert answer.status == "PYTHON_CERTIFICATE_VALID"
assert answer.certificate["cost"] == 2
assert verify(p, answer.certificate).status == "PYTHON_CERTIFICATE_VALID"
assert answer.kernel_checked is False
\end{Verbatim}

A rational entry is an integer or an exact \code{fractions.Fraction} value. Floating point entries are rejected rather than approximated. Matrix dimensions and index types are checked. The wire format represents each rational as a reduced numerator and positive denominator. Malformed input, invalid evidence, and exhausted resource budgets remain unresolved outcomes, not impossibility results.

The installed command can produce and recheck certificates. Kernel acceptance is performed by the separate source workflow with the pinned Lean kit:

\begin{Verbatim}[fontsize=\small,breaklines=true]
python scripts/certify_branch.py \
  examples/branch_corpus/Branch018_input.json \
  --kit /path/to/Komorebi_Portable_Lean_4.34.0-rc2 \
  --output /tmp/new-kernel-evidence
\end{Verbatim}

Every output directory must be new. The workflow retains the original input, the submitted or constructed certificate, the emitted Lean source, axiom lists, replay logs, and the original-input echo. The producer remains outside the formally proved portion of the system. A failed producer does not invalidate the certificate theorem, and a successful producer does not replace kernel checking.

\section{A gate for the exact distributed artifact}

A release receipt can be genuine yet stale. It may describe yesterday's source while a current wheel contains different code. A correct PDF may no longer correspond to the distributed LaTeX. A proof file may have changed after the recorded compilation. File integrity alone cannot resolve these questions unless the relevant artifacts are bound to their evidence.

The current release gate checks the strict manifest and then makes those bindings explicit. It compares every current proof source with the formal receipt, requires the complete axiom inventory, and checks that every module has a compilation and replay record. It compares the wheel's actual Python members with the source checkout and checks that installed acceptance used that same wheel. The acceptance receipt binds the Python modules, scripts, tests, README, and build metadata. A separate document receipt binds the LaTeX source and compiled PDF.

\begin{Verbatim}[fontsize=\small,breaklines=true]
python scripts/release_gate.py
\end{Verbatim}

The command is intended for a freshly extracted, unmodified release. It suppresses local bytecode creation so that checking does not itself change the sealed file inventory. Further experiments should use separate working copies or output directories.

\begin{scopebox}[title=What a passing gate means]
The gate checks consistency of the supplied source, artifacts, and receipts. It does not authenticate their author. Unsigned hashes are not signatures. An attacker able to replace the code and all receipts can fabricate mutually consistent files. The included fresh verification workflows permit the reader to perform the checks again with a separately trusted runtime.
\end{scopebox}

The gate also keeps technical readiness separate from permission to publish. The current publication decision records unselected code and document licenses. Technical acceptance is therefore reported with a pending author decision, not as authority to make a public release.

\section{Compact provenance without recursive archives}

The predecessor distribution repeatedly embedded earlier release ZIPs, which themselves embedded earlier ZIPs. The current package removes that recursive duplication without removing the active program, proof sources, examples, mathematical chapters, declaration registers, or ordinary evidence files.

Historical proof identity is retained through 121 exact source snapshots: ten from v0.2.0, 23 from v0.3.0, 37 from v0.4.0, and 51 from v0.5.0. The lineage index records their original archive hashes, filenames, byte counts, and source hashes. Each snapshot was extracted from the identified archive and compared with the current source. The full verifier repeats those source comparisons before rebuilding the current modules.

The whole predecessor ZIPs remain separate historical artifacts. They are not required to run the current verifier. Earlier reports describing bundled archives describe those earlier editions. They are retained as historical records rather than being silently rewritten into claims about this distribution.

The compiler and Mathlib cache also remain separate. The distribution contains the mathematical and software sources needed for its own workflows, not a second copy of the large pinned runtime.

\section{Acceptance and supported environment}

The formal verification rebuilds and replays all 51 modules and audits all 526 named declarations. The permitted axiom set remains \code{propext}, \code{Classical.choice}, and \code{Quot.sound}. The 24 paired rejection controls and the separate forbidden-axiom policy control are retained. The full run comprises 141 commands and 230 passing checks.

\begin{center}
\small
\begin{tabularx}{\textwidth}{@{}Xr@{}}
\toprule
Current acceptance component & Count\\
\midrule
Public API, portability, and artifact-binding unit tests & 100\\
Retained rational and menu unit tests & 134\\
Original prototype unit tests & 20\\
Original exact audit records replayed & 3,307\\
Original input/certificate pairs exercised from the wheel & 72\\
Installed-package assertions & 369\\
Installed-platform workflow checks & 24\\
\bottomrule
\end{tabularx}
\end{center}

The artifact tests include changed sources, stale wheels, absent proof coverage, forbidden axioms, duplicated records, and missing replay entries. A separate disposable-archive test refreshes the ordinary manifest after intentional corruption. The evidence bindings must still reject the altered release. These are tests of the release gate, not additional mathematical impossibility theorems.

The release-tested configuration is Linux x86\_64 with CPython 3.13.5. The C-locale regression runs with Python's UTF-8 mode disabled. File and subprocess text encodings are explicit. Paths containing spaces and Unicode are tested on that same Linux host. None of these checks constitutes a Windows simulation.

Native Windows, WSL, macOS, ARM, PyPy, and other Python interpreters remain untested. The earlier attempt to acquire additional Python versions failed at DNS resolution. A Python 3.10+ syntax or installation target is not a runtime-support claim. The offline acceptance command records whichever actual environment executes it:

\begin{Verbatim}[fontsize=\small,breaklines=true]
python scripts/accept_release.py --output /tmp/new-python-evidence
\end{Verbatim}

This creates a fresh isolated environment, installs only the local wheel, checks the retained cases, exercises the installed command, and saves its receipts. The temporary environment is removed after those receipts are written. Python acceptance is not labeled as Lean acceptance.

\section{Remaining decisions and limits}

Public code and document licenses remain the author's unmade publication decision. No license is inferred from a wish to release openly, and no GitHub, PyPI, or other public upload has been performed. A later permission or metadata change requires an appropriate artifact rebuild and renewed checks.

The external parser and source emitter are tested rather than fully formally verified. Replay uses Lean's own kernel against the pinned imports, not an independently implemented checker or a fresh rebuild of all Mathlib. Search and certificate size can be exponential. Main-library integration and authoritative Paper III imports remain separate tasks.

These limits do not require another foundational development cycle. They delimit what the standalone component proves, what its release checks establish, and which decisions remain outside the mathematics.

\backmatter
\begin{thebibliography}{14}
\bibitem{branchrelease}
Komorebi P. \emph{Constructive Repair v0.5.0: Proof Carrying Branch Search}. Structural tree checker, exact rational producer, kernel proofs, input echoes, and reproducible evidence, 19 October 2026.

\bibitem{menurelease}
Komorebi P. \emph{Constructive Repair v0.4.0: Costed Rational Menus and Rejection Covers}. Exact producer, checked cover theorems, typed certificates, verification records, and original-input echoes, 19 October 2026.

\bibitem{bridge}
Komorebi P. \emph{Constructive Repair v0.3.0: Exact Rational Certificate Bridge}. Python producer, Lean certificate modules, fixed-input records, fresh compilation and same-kernel replay evidence, 19 October 2026.

\addcontentsline{toc}{chapter}{References}
\bibitem{formal}
Komorebi P. \emph{Constructive Repair: Lean Continuation v0.2.0}. Formal source release, 19 October 2026. Ten proof modules, declaration registry, verification records, and reproducibility scripts. Historical archive: \nolinkurl{Komorebi_Constructive_Repair_Lean_v0.2.0.zip}. Exact proof snapshots and its hash are retained in the current provenance index.

\bibitem{audit}
Komorebi P. \emph{Constructive Repair Calculus: Exact Audit v0.1.0}. Mathematical audit, ordinary proofs, exact Python implementation, and case records, 19 October 2026. Preserved in the inherited material of the formal source release.

\bibitem{paperi}
Komorebi P. \emph{Projector Resolved Seed Calculus I}. Traditional monograph in Operation Komorebi. Section 11, realization export and semantic realization criterion. Source context for the preservation contract.

\bibitem{paperiii}
Komorebi P. \emph{Projector Resolved Seed Calculus III: Typed Transfer, Nonlinear Packets, and Obstruction Theory}. Modern monograph, 8 October 2026. Sections 21--24, correction maps, cokernel obstructions, nonlinear correction, and filtered lifting.

\bibitem{paperxi}
Komorebi P. \emph{Toward Arithmetic Packet Realization}. Research companion to Paper XI. Restricted constructions, failed hosts, and the remaining obligations of Conjecture 66. Source context for the distinction between finite comparison and arithmetic existence.
\bibitem{pythonio}
Python Software Foundation. \emph{Python 3.13 Library Reference: io, Text Encoding}. Documentation of locale-dependent defaults and explicit encoding selection. \url{https://docs.python.org/3.13/library/io.html#text-encoding}.
\end{thebibliography}
\end{document}
